\input{preamble}

% OK, start here.
%
\begin{document}

\title{Modules de courbes}

\maketitle

\phantomsection
\label{section-phantom}

\tableofcontents




\section{Introduction}
\label{section-introduction}

\noindent
Dans ce chapitre, nous étudions quelques-uns des champs de modules de courbes
classiques. L'article célèbre de Deligne et Mumford constitue une référence,
voir \cite{DM}.




\section{Conventions et abus de langage}
\label{section-conventions}

\noindent
Nous continuons d'employer les conventions et les abus de langage
introduits dans Propriétés des champs, section
\ref{stacks-properties-section-conventions}.
Sauf mention contraire, notre schéma de base sera $\Spec(\mathbf{Z})$.







\section{Le champ des courbes}
\label{section-stack-curves}

\noindent
Cette section fait suite à Quot, section \ref{quot-section-curves}.
Soit $\Curvesstack$ le champ dont la catégorie des sections au-dessus d'un
schéma $S$ est la catégorie des familles de courbes sur $S$.
Il importe de garder à l'esprit qu'une
{\it famille de courbes} est un morphisme $f : X \to S$, où $X$
est un espace algébrique (!) et où $f$ est plat, propre, de présentation finie
et de dimension relative $\leq 1$.
Nous savons déjà que $\Curvesstack$ est un
champ algébrique sur $\mathbf{Z}$, voir Quot, théorème
\ref{quot-theorem-curves-algebraic}.
Si nous n'autorisions pas les espaces algébriques dans la définition de
notre champ, ce théorème serait faux.

\medskip\noindent
On note souvent un changement de base par un indice, mais nous ne pouvons pas
employer cette notation pour $\Curvesstack$, car $\Curvesstack_S$
désigne chez nous la catégorie fibre au-dessus de $S$.
C'est pourquoi, dans Quot, remarque \ref{quot-remark-curves-base-change},
nous avons noté $B\text{-}\Curvesstack$ le changement de base
$$
B\text{-}\Curvesstack = \Curvesstack \times B
$$
à l'espace algébrique $B$. Le produit de droite est pris au-dessus de
l'objet final, c'est-à-dire au-dessus de $\Spec(\mathbf{Z})$. L'objet de gauche
est le champ classifiant les familles de courbes sur la catégorie des schémas
sur $B$. En particulier, si $k$ est un corps, alors
$$
k\text{-}\Curvesstack = \Curvesstack \times \Spec(k)
$$
est le champ de modules classifiant les familles de courbes sur la catégorie
des schémas sur $k$.
Avant de poursuivre, effectuons une vérification de cohérence.

\begin{lemma}
\label{lemma-extend-curves-to-spaces}
Soit $T \to B$ un morphisme d'espaces algébriques. La catégorie
$$
\Mor_B(T, B\text{-}\Curvesstack) = \Mor(T, \Curvesstack)
$$
est la catégorie des familles de courbes sur $T$.
\end{lemma}

\begin{proof}
Une famille de courbes sur $T$ est un morphisme $f : X \to T$ d'espaces
algébriques, qui est plat, propre, de présentation finie et de
dimension relative $\leq 1$ (Morphismes d'espaces, définition
\ref{spaces-morphisms-definition-relative-dimension}).
C'est exactement la même définition que dans Quot, situation
\ref{quot-situation-curves}, si ce n'est que la base $T$ peut être
un espace algébrique. Notre catégorie de base par défaut pour les
champs et espaces algébriques est la catégorie des schémas ; le lemme
ne découle donc pas immédiatement des définitions. Cela étant dit,
nous invitons le lecteur à sauter la démonstration.

\medskip\noindent
D'après la description de $B\text{-}\Curvesstack$ comme produit donnée
ci-dessus, il suffit de démontrer le lemme dans le cas absolu. Choisissons
un schéma $U$ et un morphisme étale surjectif $p : U \to T$.
Posons $R = U \times_T U$, avec projections $s, t : R \to U$.

\medskip\noindent
Soit $v : T \to \Curvesstack$ un morphisme. Alors $v \circ p$
correspond à une famille de courbes $X_U \to U$. Le
$2$-morphisme canonique $v \circ p \circ t \to v \circ p \circ s$
est un isomorphisme $\varphi : X_U \times_{U, s} R \to X_U \times_{U, t} R$.
Cet isomorphisme satisfait à la condition de cocycle sur
$R \times_{s, t} R$.
D'après Amorçage, lemme \ref{bootstrap-lemma-descend-algebraic-space},
nous obtenons un morphisme d'espaces algébriques $X \to T$
dont l'image inverse sur $U$ est égale à $X_U$, compatiblement avec $\varphi$.
Puisque $\{U \to T\}$ est un recouvrement étale, le morphisme
$X \to T$ est plat, propre et de présentation finie d'après
Descente pour les espaces, lemmes
\ref{spaces-descent-lemma-descending-property-flat},
\ref{spaces-descent-lemma-descending-property-proper} et
\ref{spaces-descent-lemma-descending-property-finite-presentation}.
De plus, $X \to T$ est de dimension relative $\leq 1$, car cette propriété
est locale pour la topologie étale. Ainsi, $X \to T$ est une famille de
courbes sur $T$.

\medskip\noindent
Réciproquement, soit $X \to T$ une famille de courbes. Alors le
changement de base $X_U$ détermine un morphisme $w : U \to \Curvesstack$,
et l'isomorphisme canonique $X_U \times_{U, s} R \to X_U \times_{U, t} R$
détermine une $2$-flèche $w \circ s \to w \circ t$ satisfaisant à la
condition de cocycle. On obtient ainsi un morphisme
$v : T = [U/R] \to \Curvesstack$ par la propriété universelle du
quotient $[U/R]$ ; voir
Groupoïdes dans les espaces, lemme
\ref{spaces-groupoids-lemma-quotient-stack-2-coequalizer}.
(En fait, il est ici bien plus simple de revenir à la situation antérieure
à l'introduction de notre abus de langage et de construire directement
le foncteur $\Sch/T \to \Curvesstack$ qui « est » le
morphisme $T \to \Curvesstack$.)

\medskip\noindent
Nous omettons de vérifier que les constructions ci-dessus
s'étendent aux morphismes entre objets et sont quasi-inverses l'une de l'autre.
\end{proof}






\section{Le champ des courbes polarisées}
\label{section-polarized-curves}

\noindent
Dans cette section, nous développons une partie des considérations de
Quot, remarque \ref{quot-remark-alternative-approach-curves}.
Considérons le $2$-produit fibré
$$
\xymatrix{
\Curvesstack \times_{\Spacesstack'_{fp, flat, proper}}
\Polarizedstack \ar[r] \ar[d] &
\Polarizedstack \ar[d] \\
\Curvesstack \ar[r] &
\Spacesstack'_{fp, flat, proper}
}
$$
Nous notons ce $2$-produit fibré
$$
\textit{PolarizedCurves} =
\Curvesstack
\times_{\Spacesstack'_{fp, flat, proper}}
\Polarizedstack
$$
Ce produit fibré paramètre les courbes polarisées, c'est-à-dire les familles
de courbes munies d'un faisceau inversible relativement ample.
Plus précisément, un objet de
$\textit{PolarizedCurves}$
est un couple $(X \to S, \mathcal{L})$ tel que
\begin{enumerate}
\item $X \to S$ est un morphisme de schémas propre, plat,
de présentation finie et de dimension relative $\leq 1$, et
\item $\mathcal{L}$ est un $\mathcal{O}_X$-module inversible
qui est relativement ample sur $X/S$.
\end{enumerate}
Un morphisme $(X' \to S', \mathcal{L}') \to (X \to S, \mathcal{L})$
entre objets de
$\textit{PolarizedCurves}$
est donné par un triplet $(f, g, \varphi)$,
où $f : X' \to X$ et $g : S' \to S$
sont des morphismes de schémas s'inscrivant dans un diagramme commutatif
$$
\xymatrix{
X' \ar[d] \ar[r]_f & X \ar[d] \\
S' \ar[r]^g & S
}
$$
qui induit un isomorphisme $X' \to S' \times_S X$ ; autrement dit, le
diagramme est cartésien, et $\varphi : f^*\mathcal{L} \to \mathcal{L}'$
est un isomorphisme. La composition se définit de manière évidente.

\begin{lemma}
\label{lemma-polarized-curves-in-polarized}
Le morphisme
$\textit{PolarizedCurves} \to
\Polarizedstack$ est une immersion ouverte et fermée.
\end{lemma}

\begin{proof}
Cela résulte de ce que le $1$-morphisme
$\Curvesstack \to \Spacesstack'_{fp, flat, proper}$
est représentable par des immersions ouvertes et fermées ; voir
Quot, lemme \ref{quot-lemma-curves-open-and-closed-in-spaces}.
\end{proof}

\begin{lemma}
\label{lemma-polarized-curves-over-curves}
Le morphisme
$\textit{PolarizedCurves} \to \Curvesstack$
est lisse et surjectif.
\end{lemma}

\begin{proof}
Surjectivité. Soient $k$ un corps et $X$ un espace algébrique propre
sur $k$, de dimension $\leq 1$, c'est-à-dire un objet de $\Curvesstack$ sur $k$.
D'après Espaces sur les corps, lemme
\ref{spaces-over-fields-lemma-codim-1-point-in-schematic-locus}
l'espace algébrique $X$ est un schéma. Ainsi, $X$
est un schéma propre de dimension $\leq 1$ sur $k$.
D'après Variétés, lemme \ref{varieties-lemma-dim-1-proper-projective},
$X$ est H-projectif sur $\kappa$.
Il existe en particulier un $\mathcal{O}_X$-module inversible ample
$\mathcal{L}$ sur $X$. Alors $(X, \mathcal{L})$ est un objet
de $\textit{PolarizedCurves}$ sur
$k$ dont l'image est $X$.

\medskip\noindent
Lissité. Soit $X \to S$ un objet de $\Curvesstack$, c'est-à-dire un
morphisme $S \to \Curvesstack$. Il est clair que
$$
\textit{PolarizedCurves}
\times_{\Curvesstack} S
\subset \Picardstack_{X/S}
$$
est le sous-champ des objets $(T/S, \mathcal{L}/X_T)$ tels que
$\mathcal{L}$ soit ample sur $X_T/T$. C'est un sous-champ ouvert d'après
Descente pour les espaces, lemme \ref{spaces-descent-lemma-ample-in-neighbourhood}.
Puisque $\Picardstack_{X/S} \to S$ est lisse d'après
Champs de modules, lemme \ref{moduli-lemma-pic-curves-smooth},
la conclusion s'ensuit.
\end{proof}

\begin{lemma}
\label{lemma-etale-locally-scheme}
Soit $X \to S$ une famille de courbes.
Il existe alors un recouvrement étale $\{S_i \to S\}$
tel que $X_i = X \times_S S_i$ soit un schéma. Nous pouvons même
supposer $X_i$ H-projectif sur $S_i$.
\end{lemma}

\begin{proof}
Il s'agit d'un corollaire immédiat du
lemme \ref{lemma-polarized-curves-over-curves}.
En effet, l'explicitation des définitions montre que ce lemme fournit un
morphisme lisse surjectif $S' \to S$ tel que $X' = X \times_S S'$
soit muni d'un $\mathcal{O}_{X'}$-module inversible
$\mathcal{L}'$ qui est ample sur $X'/S'$.
Nous pouvons ensuite raffiner le recouvrement lisse $\{S' \to S\}$
par un recouvrement étale $\{S_i \to S\}$ ; voir
Compléments sur les morphismes, lemme \ref{more-morphisms-lemma-etale-dominates-smooth}.
Après remplacement de $S_i$ par un recouvrement ouvert convenable, nous
pouvons supposer $X_i \to S_i$ H-projectif ; voir
Morphismes, lemmes \ref{morphisms-lemma-proper-ample-locally-projective} et
\ref{morphisms-lemma-characterize-locally-projective}
(ce point est aussi étudié en détail dans
Compléments sur les morphismes, section \ref{more-morphisms-section-projective}).
\end{proof}






\section{Propriétés du champ des courbes}
\label{section-properties}

\noindent
Le lemme suivant n'est pas vrai pour les modules de surfaces ; voir
la remarque \ref{remark-boundedness-aut-does-not-work-surfaces}.

\begin{lemma}
\label{lemma-curves-diagonal-separated-fp}
La diagonale de $\Curvesstack$ est séparée
et de présentation finie.
\end{lemma}

\begin{proof}
Rappelons que $\Curvesstack$ est un champ algébrique qui préserve les limites ;
voir Quot, lemme \ref{quot-lemma-curves-limits}.
D'après Limites de champs, lemme
\ref{stacks-limits-lemma-limit-preserving-diagonal}, cela implique que
$\Delta : \Polarizedstack \to \Polarizedstack \times \Polarizedstack$
préserve les limites. Par conséquent, $\Delta$ est localement de présentation
finie d'après Limites de champs, proposition
\ref{stacks-limits-proposition-characterize-locally-finite-presentation}.

\medskip\noindent
Montrons que $\Delta$ est séparée. Pour cela, il suffit de montrer que,
pour un schéma $U$ et deux objets $Y \to U$ et $X \to U$ de
$\Curvesstack$ au-dessus de $U$, l'espace algébrique
$$
\mathit{Isom}_U(Y, X)
$$
est séparé. Il résulte en effet de Champs de modules, lemmes
\ref{moduli-lemma-Mor-s-lfp} et
\ref{moduli-lemma-Isom-in-Mor}, que cet espace
algébrique est séparé.

\medskip\noindent
Pour achever la démonstration, montrons que $\Delta$ est quasi-compacte.
Puisque $\Delta$ est représentable par des espaces algébriques, il suffit
de vérifier que le changement de base de $\Delta$ par un morphisme lisse
surjectif $U \to \Curvesstack \times \Curvesstack$ est quasi-compact
(voir, par exemple, Propriétés des champs, lemme
\ref{stacks-properties-lemma-check-property-covering}).
Choisissons pour $U = \coprod U_i$ une réunion disjointe d'ouverts affines,
munie d'un morphisme lisse surjectif
$$
U \longrightarrow
\textit{PolarizedCurves} \times \textit{PolarizedCurves}
$$
Alors $U \to \Curvesstack \times \Curvesstack$ est surjectif
et lisse, car $\textit{PolarizedCurves} \to \Curvesstack$
est surjectif et lisse d'après le lemme
\ref{lemma-polarized-curves-over-curves}.
Comme $\textit{PolarizedCurves}$ préserve les limites
(d'après Axiomes d'Artin, lemme
\ref{artin-lemma-fibre-product-limit-preserving},
et Quot, lemmes \ref{quot-lemma-curves-limits},
\ref{quot-lemma-polarized-limits} et
\ref{quot-lemma-spaces-limits}), le morphisme
$\textit{PolarizedCurves} \to \Spec(\mathbf{Z})$ est localement de
présentation finie. Il en va donc de même de $U_i \to \Spec(\mathbf{Z})$
(Limites de champs, proposition
\ref{stacks-limits-proposition-characterize-locally-finite-presentation}
et Morphismes de champs, lemmes
\ref{stacks-morphisms-lemma-composition-finite-presentation} et
\ref{stacks-morphisms-lemma-smooth-locally-finite-presentation}).
En particulier, $U_i$ est affine noethérien. Nous sommes ainsi ramenés
au cas étudié dans le paragraphe suivant.

\medskip\noindent
Dans ce paragraphe, étant donnés un schéma affine noethérien $U$ et deux objets
$(Y, \mathcal{N})$ et $(X, \mathcal{L})$
de $\textit{PolarizedCurves}$ au-dessus de $U$, nous montrons que l'espace algébrique
$$
\mathit{Isom}_U(Y, X)
$$
est quasi-compact. Comme les composantes connexes de $U$ sont ouvertes et
fermées, nous pouvons remplacer $U$ par celles-ci. Nous supposons donc
désormais $U$ connexe. Soit $u \in U$ un point. Soient $Q$ et $P$ les
polynômes de Hilbert de ces familles, c'est-à-dire
$$
Q(n) = \chi(Y_u, \mathcal{N}_u^{\otimes n})
\quad\text{et}\quad
P(n) = \chi(X_u, \mathcal{L}_u^{\otimes n})
$$
voir Variétés, lemme \ref{varieties-lemma-numerical-polynomial-from-euler}.
Comme $U$ est connexe et que les fonctions
$u \mapsto \chi(Y_u, \mathcal{N}_u^{\otimes n})$ et
$u \mapsto \chi(X_u, \mathcal{L}_u^{\otimes n})$
 sont localement constantes (voir
Catégories dérivées de schémas, lemme
\ref{perfect-lemma-chi-locally-constant-geometric})
nous obtenons le même polynôme de Hilbert en tout point de $U$.
Posons
$$
\mathcal{M} = \text{pr}_1^*\mathcal{N}
\otimes_{\mathcal{O}_{Y \times_U X}} \text{pr}_2^*\mathcal{L}
$$
sur $Y \times_U X$. Pour
$(f, \varphi) \in \mathit{Isom}_U(Y, X)(T)$, où $T$ est un schéma sur $U$,
et pour tout $t \in T$, on a
\begin{align*}
\chi(Y_t, (\text{id} \times f)^*\mathcal{M}^{\otimes n})
& =
\chi(Y_t,
\mathcal{N}_t^{\otimes n} \otimes_{\mathcal{O}_{Y_t}}
f_t^*\mathcal{L}_t^{\otimes n}) \\
& =
n\deg(\mathcal{N}_t) + n\deg(f_t^*\mathcal{L}_t) +
\chi(Y_t, \mathcal{O}_{Y_t}) \\
& =
Q(n) + n\deg(\mathcal{L}_t) \\
& =
Q(n) + P(n) - P(0)
\end{align*}
d'après Riemann--Roch pour les courbes propres, plus précisément d'après
Variétés, définition \ref{varieties-definition-degree-invertible-sheaf},
lemme \ref{varieties-lemma-degree-tensor-product},
et le fait que $f_t$ est un isomorphisme.
En posant $P'(t) = Q(t) + P(t) - P(0)$, nous obtenons
$$
\mathit{Isom}_U(Y, X) =
\mathit{Isom}_U(Y, X) \cap \mathit{Mor}^{P', \mathcal{M}}_U(Y, X)
$$
Cette intersection est une intersection de sous-espaces ouverts de
$\mathit{Mor}_U(Y, X)$ ; voir Champs de modules, lemme
\ref{moduli-lemma-Isom-in-Mor} et remarque
\ref{moduli-remark-Mor-numerical}.
Or $\mathit{Mor}^{P', \mathcal{M}}_U(Y, X)$
est un espace algébrique noethérien, puisqu'il est de présentation
finie sur $U$ d'après Champs de modules, lemme
\ref{moduli-lemma-Mor-qc-over-base}.
L'intersection est donc elle aussi un espace algébrique noethérien,
ce qui achève la démonstration.
\end{proof}

\begin{remark}
\label{remark-boundedness-aut-does-not-work-surfaces}
L'argument de bornitude employé dans la démonstration du
lemme \ref{lemma-curves-diagonal-separated-fp}
ne fonctionne pas pour les modules de surfaces ; de fait,
le résultat est faux, notamment parce que des surfaces K3
sur des corps peuvent avoir des groupes discrets d'automorphismes infinis.
La « raison » pour laquelle l'argument ne fonctionne pas est que, sur une
surface projective $S$ sur un corps, étant donnés des faisceaux inversibles
amples $\mathcal{N}$ et $\mathcal{L}$ de polynômes de Hilbert $Q$ et $P$,
il n'existe aucune borne a priori pour le polynôme de Hilbert de
$\mathcal{N} \otimes_{\mathcal{O}_S} \mathcal{L}$.
En termes de théorie de l'intersection, si $H_1$ et $H_2$ sont des
diviseurs de Cartier effectifs et amples sur $S$, il n'existe aucune
borne (supérieure) pour le nombre d'intersection $H_1 \cdot H_2$
en fonction de $H_1 \cdot H_1$ et de $H_2 \cdot H_2$.
\end{remark}

\begin{lemma}
\label{lemma-curves-qs-lfp}
Le morphisme $\Curvesstack \to \Spec(\mathbf{Z})$ est quasi-séparé et
localement de présentation finie.
\end{lemma}

\begin{proof}
Pour vérifier que $\Curvesstack \to \Spec(\mathbf{Z})$ est quasi-séparé,
il faut montrer que sa diagonale est quasi-compacte et quasi-séparée.
Cela résulte immédiatement du lemme
\ref{lemma-curves-diagonal-separated-fp}.
Pour montrer que $\Curvesstack \to \Spec(\mathbf{Z})$ est localement de
présentation finie, il suffit de montrer que $\Curvesstack$
préserve les limites ; voir Limites de champs, proposition
\ref{stacks-limits-proposition-characterize-locally-finite-presentation}.
C'est Quot, lemme \ref{quot-lemma-curves-limits}.
\end{proof}






\section{Sous-champs ouverts du champ des courbes}
\label{section-open}

\noindent
Dans la suite, nous caractériserons souvent un sous-champ ouvert de
$\Curvesstack$ au moyen d'une propriété $P$ des morphismes d'espaces
algébriques. Pour voir que $P$ définit un sous-champ ouvert, il suffit
de vérifier la condition suivante :
\begin{enumerate}
\item[(o)] pour toute famille de courbes $f : X \to S$, il existe
un plus grand sous-schéma ouvert $S' \subset S$ tel que
$f|_{f^{-1}(S')} : f^{-1}(S') \to S'$ possède $P$ et tel que
la formation de $S'$ commute à tout changement de base.
\end{enumerate}
Supposons en effet que (o) soit satisfaite. Choisissons un schéma $U$ et
un morphisme lisse surjectif $m : U \to \Curvesstack$. Posons
$R = U \times_{\Curvesstack} U$ et notons $t, s : R \to U$
les projections. Rappelons que $\Curvesstack = [U/R]$ est une présentation ;
voir Champs algébriques, lemme \ref{algebraic-lemma-stack-presentation} et
définition \ref{algebraic-definition-presentation}.
Par construction de $\Curvesstack$ comme champ des courbes, le morphisme $m$
est le morphisme classifiant d'une famille de courbes $C \to U$. La
$2$-commutativité du diagramme
$$
\xymatrix{
R \ar[r]_s \ar[d]_t & U \ar[d] \\
U \ar[r] & \Curvesstack
}
$$
implique que $C \times_{U, s} R  \cong C \times_{U, t} R$
(isomorphisme de familles de courbes sur $R$). Soit $W \subset U$
le plus grand sous-schéma ouvert tel que
$f|_{f^{-1}(W)} : f^{-1}(W) \to W$ possède $P$, comme dans (o).
Puisque, d'après (o), la formation de $W$ commute aux changements de base,
l'isomorphisme ci-dessus donne $s^{-1}(W) = t^{-1}(W)$.
Ainsi, $W \subset U$ correspond à un sous-champ ouvert
$$
\Curvesstack^P \subset \Curvesstack
$$
d'après Propriétés des champs, lemme
\ref{stacks-properties-lemma-immersion-presentation}.

\medskip\noindent
En conservant les notations du paragraphe précédent, nous affirmons que
le sous-champ ouvert $\Curvesstack^P$ possède les deux propriétés
universelles suivantes :
\begin{enumerate}
\item pour une famille de courbes $X \to S$, les assertions suivantes
sont équivalentes :
\begin{enumerate}
\item le morphisme classifiant $S \to \Curvesstack$ se factorise par
$\Curvesstack^P$,
\item le morphisme $X \to S$ possède $P$ ;
\end{enumerate}
\item pour un schéma $X$ propre sur un corps $k$, de dimension $\leq 1$,
les assertions suivantes sont équivalentes :
\begin{enumerate}
\item le morphisme classifiant $\Spec(k) \to \Curvesstack$ se factorise
par $\Curvesstack^P$,
\item le morphisme $X \to \Spec(k)$ possède $P$.
\end{enumerate}
\end{enumerate}
Cela résulte de la considération du $2$-produit fibré
$$
\xymatrix{
T \ar[r]_p \ar[d]_q & U \ar[d] \\
S \ar[r] & \Curvesstack
}
$$
Observons que $T \to S$ est surjectif et lisse, car il se déduit de
$U \to \Curvesstack$ par changement de base. L'ouvert $S' \subset S$
donné par (o) est donc déterminé par son image inverse dans $T$.
Or, par l'invariance des ouverts de (o) sous changement de base et puisque
$X \times_S T \cong C \times_U T$ par $2$-commutativité, nous obtenons
$q^{-1}(S') = p^{-1}(W)$ comme ouverts de $T$.
Cela implique immédiatement (1). La partie (2) est un cas particulier de (1).

\medskip\noindent
Étant données deux propriétés $P$ et $Q$ des morphismes d'espaces
algébriques, supposons que nous ayons déjà établi que $\Curvesstack^Q$
est un sous-champ ouvert de $\Curvesstack$. La même méthode permet alors
de démontrer que
$\Curvesstack^{Q, P} \subset \Curvesstack^Q$ est ouvert.
Nous omettons les détails.



\section{Courbes à groupes d'automorphismes finis et réduits}
\label{section-finite-aut}

\noindent
Soit $X$ un schéma propre sur un corps $k$, de dimension $\leq 1$,
c'est-à-dire un objet de $\Curvesstack$ sur $k$.
D'après le lemme \ref{lemma-curves-diagonal-separated-fp},
l'espace algébrique des automorphismes $\mathit{Aut}(X)$
est de type fini et séparé sur $k$.
En particulier, $\mathit{Aut}(X)$ est un schéma en groupes ; voir
Compléments sur les groupoïdes dans les espaces, lemme
\ref{spaces-more-groupoids-lemma-group-space-scheme-locally-finite-type-over-k}.
Si la caractéristique de $k$ est nulle, alors $\mathit{Aut}(X)$
est réduit et même lisse sur $k$ (Groupoïdes, lemme
\ref{groupoids-lemma-group-scheme-characteristic-zero-smooth}).
Toutefois, en général, $\mathit{Aut}(X)$ n'est pas réduit, même
si $X$ est géométriquement réduit.

\begin{example}[Groupe d'automorphismes non réduit]
\label{example-non-reduced}
Soit $k$ un corps algébriquement clos de caractéristique $2$.
Posons $Y = Z = \mathbf{P}^1_k$. Choisissons trois points $k$-rationnels
deux à deux distincts $a, b, c$ dans $\mathbf{A}^1_k$. En considérant
$\mathbf{A}^1_k \subset \mathbf{P}^1_k = Y = Z$ comme un sous-schéma ouvert,
nous obtenons une immersion fermée
$$
T =  \Spec(k[t]/(t - a)^2) \amalg \Spec(k[t]/(t - b)^2)
\amalg \Spec(k[t]/(t - c)^2)
\longrightarrow
\mathbf{P}^1_k
$$
Soit $X$ la somme amalgamée dans le diagramme
$$
\xymatrix{
T \ar[r] \ar[d] & Y \ar[d] \\
Z \ar[r] & X
}
$$
Soit $U \subset X$ la partie ouverte affine qui est l'image de
$\mathbf{A}^1_k \amalg \mathbf{A}^1_k$. Nous avons alors un diagramme
égalisateur
$$
\xymatrix{
\mathcal{O}_X(U) \ar[r] &
k[t] \times k[t] \ar@<1ex>[r] \ar@<-1ex>[r] &
k[t]/(t - a)^2 \times k[t]/(t - b)^2 \times k[t]/(t - c)^2
}
$$
Sur l'algèbre des nombres duaux $A = k[\epsilon]$, ce diagramme égalisateur
admet un automorphisme non trivial envoyant $t$ sur $t + \epsilon$. Nous
laissons au lecteur le soin de voir que cet automorphisme s'étend en un
automorphisme de $X$ sur $A$. D'autre part, le lecteur vérifie aisément que
le groupe d'automorphismes de $X$ sur $k$ est fini.
Ainsi, $\mathit{Aut}(X)$ est nécessairement non réduit.
\end{example}

\noindent
Soit $X$ un schéma propre sur un corps $k$, de dimension $\leq 1$,
c'est-à-dire un objet de $\Curvesstack$ sur $k$. Même si $\mathit{Aut}(X)$
est géométriquement réduit, il n'est pas nécessairement de
dimension $0$, quand bien même $X$ est lisse et géométriquement connexe.

\begin{example}[Groupe d'automorphismes lisse de dimension positive]
\label{example-pos-dim}
Soit $k$ un corps algébriquement clos. Si $X$ est une courbe lisse
de genre $0$, resp.\ $1$, alors son groupe d'automorphismes est de
dimension $3$, resp.\ $1$. En effet, dans le cas du genre $0$, on a
$X \cong \mathbf{P}^1_k$ d'après Courbes algébriques, proposition
\ref{curves-proposition-projective-line}. Puisque
$$
\mathit{Aut}(\mathbf{P}^1_k) = \text{PGL}_{2, k}
$$
comme foncteurs, sa dimension est $3$. D'autre part,
si le genre de $X$ est $1$, alors l'application
$X = \underline{\Hilbfunctor}^1_{X/k} \to
\underline{\Picardfunctor}^1_{X/k}$ est un isomorphisme ; voir
Schémas de Picard des courbes, lemme \ref{pic-lemma-picard-pieces},
et Courbes algébriques, théorème
\ref{curves-theorem-curves-rational-maps}.
Ainsi, $X$ possède une structure de variété abélienne
(puisque $\underline{\Picardfunctor}^1_{X/k} \cong
\underline{\Picardfunctor}^0_{X/k}$).
En particulier, les fibrés tangent et cotangent de $X$ sont triviaux
(Groupoïdes, lemme \ref{groupoids-lemma-group-scheme-module-differentials}).
Nous en concluons que $\dim_k H^0(X, T_X) = 1$, donc que
$\dim \mathit{Aut}(X) \leq 1$. D'autre part, les translations
(en considérant $X$ comme un schéma en groupes) fournissent une partie
de dimension $1$ de $\text{Aut}(X)$ ; sa dimension vaut donc bien $1$.
\end{example}

\noindent
Il existe en fait un sous-champ ouvert de
$\Curvesstack$ qui paramètre les courbes dont le groupe d'automorphismes
est géométriquement réduit et fini.
En voici un énoncé précis.

\begin{lemma}
\label{lemma-DM-curves}
Il existe un sous-champ ouvert $\Curvesstack^{DM} \subset \Curvesstack$
possédant les propriétés suivantes :
\begin{enumerate}
\item $\Curvesstack^{DM} \subset \Curvesstack$ est le plus grand
sous-champ ouvert qui soit DM ;
\item pour une famille de courbes $X \to S$, les assertions suivantes
sont équivalentes :
\begin{enumerate}
\item le morphisme classifiant $S \to \Curvesstack$ se factorise par
$\Curvesstack^{DM}$,
\item l'espace algébrique en groupes $\mathit{Aut}_S(X)$ est non ramifié
sur $S$ ;
\end{enumerate}
\item pour un schéma $X$ propre sur un corps $k$, de dimension $\leq 1$,
les assertions suivantes sont équivalentes :
\begin{enumerate}
\item le morphisme classifiant $\Spec(k) \to \Curvesstack$ se factorise
par $\Curvesstack^{DM}$ ;
\item $\mathit{Aut}(X)$ est géométriquement réduit sur $k$ et
de dimension $0$ ;
\item $\mathit{Aut}(X) \to \Spec(k)$ est non ramifié.
\end{enumerate}
\end{enumerate}
\end{lemma}

\begin{proof}
L'existence d'un sous-champ ouvert ayant la propriété (1) résulte de
Morphismes de champs, lemme \ref{stacks-morphisms-lemma-open-DM-locus}.
Les points de ce sous-champ ouvert sont caractérisés par (3)(c) d'après
Morphismes de champs, lemme \ref{stacks-morphisms-lemma-points-DM-locus}.
L'équivalence de (3)(b) et (3)(c) affirme qu'un espace algébrique $G$,
localement de type fini, géométriquement réduit et de dimension $0$ sur
un corps $k$, est non ramifié sur $k$.
Tout d'abord, $G$ est un schéma d'après Espaces sur les corps, lemme
\ref{spaces-over-fields-lemma-locally-finite-type-dim-zero}.
Nous pouvons ensuite prendre un ouvert affine de $G$, observer
qu'il est propre sur $k$, puis appliquer Variétés, lemme
\ref{varieties-lemma-proper-geometrically-reduced-global-sections}.
Nous omettons les détails mineurs.

\medskip\noindent
La partie (2) est vraie parce que (3) l'est. En effet, le morphisme
$\mathit{Aut}_S(X) \to S$ est localement de type fini. Nous pouvons donc
vérifier si $\mathit{Aut}_S(X) \to S$ est non ramifié en tout point de
$\mathit{Aut}_S(X)$ en le vérifiant sur les fibres aux points du schéma $S$ ;
voir Morphismes d'espaces, lemme
\ref{spaces-morphisms-lemma-where-unramified}.
Mais, après changement de base en un point de $S$, nous nous ramenons à
l'équivalence de (3)(a) et (3)(c).
\end{proof}

\begin{lemma}
\label{lemma-in-DM-locus-vector-fields}
Soit $X$ un schéma propre sur un corps $k$, de dimension $\leq 1$.
Alors les propriétés (3)(a), (b), (c) sont également équivalentes à
$\text{Der}_k(\mathcal{O}_X, \mathcal{O}_X) = 0$.
\end{lemma}

\begin{proof}
Dans la discussion précédente, nous avons vu que $G = \mathit{Aut}(X)$
est un schéma en groupes de type fini et séparé sur $\Spec(k)$ ;
on utilise ici le lemme \ref{lemma-curves-diagonal-separated-fp} et
Compléments sur les groupoïdes dans les espaces, lemme
\ref{spaces-more-groupoids-lemma-group-space-scheme-locally-finite-type-over-k}.
Alors $G$ est non ramifié sur $k$ si et seulement si $\Omega_{G/k} = 0$
(Morphismes, lemme \ref{morphisms-lemma-unramified-omega-zero}).
D'après Groupoïdes, lemme
\ref{groupoids-lemma-group-scheme-module-differentials}, cette annulation
a lieu si $T_{G/k, e} = 0$, où $T_{G/k, e}$ est l'espace tangent à $G$
en l'élément neutre $e \in G(k)$ ; voir Variétés, définition
\ref{varieties-definition-tangent-space}, ainsi que la formule de
Variétés, lemme \ref{varieties-lemma-tangent-space-cotangent-space}.
Puisque $\kappa(e) = k$, l'espace tangent se décrit au moyen des
morphismes $\alpha : \Spec(k[\epsilon]) \to G = \mathit{Aut}(X)$
dont la restriction à $\Spec(k)$ est $e$.
On est ainsi ramené à considérer les automorphismes
$$
\alpha :
X \times_{\Spec(k)} \Spec(k[\epsilon])
\longrightarrow
X \times_{\Spec(k)} \Spec(k[\epsilon])
$$
sur $\Spec(k[\epsilon])$ dont la restriction à $\Spec(k)$ est
$\text{id}_X$. De tels automorphismes sont appelés
automorphismes infinitésimaux.

\medskip\noindent
Les automorphismes infinitésimaux de $X$ correspondent bijectivement
aux dérivations de $\mathcal{O}_X$ sur $k$. Cela résulte de
Compléments sur les morphismes, lemmes
\ref{more-morphisms-lemma-difference-derivation} et
\ref{more-morphisms-lemma-action-by-derivations} (nous n'avons besoin que
du premier, puisque le sens réciproque ne nous intéresse pas ; consulter
aussi Compléments sur les morphismes, remarque
\ref{more-morphisms-remark-another-special-case}, pour une explication).
Pour un autre argument établissant cette égalité, nous renvoyons le lecteur
à Problèmes de déformation, lemme \ref{examples-defos-lemma-schemes-TI}.
\end{proof}





\section{Courbes de Cohen-Macaulay}
\label{section-CM}

\noindent
Il existe un sous-champ ouvert de $\Curvesstack$ qui paramètre
les « courbes » de Cohen-Macaulay.

\begin{lemma}
\label{lemma-CM-curves}
Il existe un sous-champ ouvert $\Curvesstack^{CM} \subset \Curvesstack$
tel que :
\begin{enumerate}
\item pour une famille de courbes $X \to S$, les assertions suivantes
sont équivalentes :
\begin{enumerate}
\item le morphisme classifiant $S \to \Curvesstack$ se factorise par
$\Curvesstack^{CM}$,
\item le morphisme $X \to S$ est Cohen-Macaulay ;
\end{enumerate}
\item pour un schéma $X$ propre sur un corps $k$ et tel que $\dim(X) \leq 1$,
les assertions suivantes sont équivalentes :
\begin{enumerate}
\item le morphisme classifiant $\Spec(k) \to \Curvesstack$ se factorise
par $\Curvesstack^{CM}$,
\item $X$ est Cohen-Macaulay.
\end{enumerate}
\end{enumerate}
\end{lemma}

\begin{proof}
Soit $f : X \to S$ une famille de courbes. D'après
Compléments sur les morphismes d'espaces, lemme
\ref{spaces-more-morphisms-lemma-flat-finite-presentation-CM-open},
l'ensemble
$$
W = \{x \in |X| : f \text{ est Cohen-Macaulay au point }x\}
$$
est ouvert dans $|X|$, et la formation de cet ouvert commute à tout
changement de base. Puisque $f$ est propre, le sous-ensemble
$$
S' = S \setminus f(|X| \setminus W)
$$
de $S$ est ouvert, et $X \times_S S' \to S'$ est Cohen-Macaulay.
De plus, la formation de $S'$ commute à tout changement de base,
car il en est ainsi de $W$. La méthode exposée à la section
\ref{section-open} fournit donc le sous-champ ouvert possédant les
propriétés voulues.
\end{proof}

\begin{lemma}
\label{lemma-CM-1-curves}
Il existe un sous-champ ouvert $\Curvesstack^{CM, 1} \subset \Curvesstack$
tel que :
\begin{enumerate}
\item pour une famille de courbes $X \to S$, les assertions suivantes
sont équivalentes :
\begin{enumerate}
\item le morphisme classifiant $S \to \Curvesstack$ se factorise par
$\Curvesstack^{CM, 1}$,
\item le morphisme $X \to S$ est Cohen-Macaulay et de dimension
relative $1$ (Morphismes d'espaces, définition
\ref{spaces-morphisms-definition-relative-dimension}) ;
\end{enumerate}
\item pour un schéma $X$ propre sur un corps $k$ et tel que $\dim(X) \leq 1$,
les assertions suivantes sont équivalentes :
\begin{enumerate}
\item le morphisme classifiant $\Spec(k) \to \Curvesstack$ se factorise
par $\Curvesstack^{CM, 1}$,
\item $X$ est Cohen-Macaulay et $X$ est équidimensionnel de dimension $1$.
\end{enumerate}
\end{enumerate}
\end{lemma}

\begin{proof}
D'après le lemme \ref{lemma-CM-curves}, il est clair que l'on a
$\Curvesstack^{CM, 1} \subset \Curvesstack^{CM}$ si ce sous-champ existe.
Soit $f : X \to S$ une famille de courbes telle que $f$ soit
Cohen-Macaulay. D'après Compléments sur les morphismes d'espaces, lemme
\ref{spaces-more-morphisms-lemma-lfp-CM-relative-dimension},
nous avons une décomposition
$$
X = X_0 \amalg X_1
$$
en sous-espaces ouverts et fermés telle que $X_0 \to S$ soit de dimension
relative $0$ et $X_1 \to S$ de dimension relative $1$.
Puisque $f$ est propre, le sous-ensemble
$$
S' = S \setminus f(|X_0|)
$$
de $S$ est ouvert, et $X \times_S S' \to S'$ est Cohen-Macaulay
et de dimension relative $1$.
De plus, la formation de $S'$ commute à tout changement de base,
car il en est ainsi de la décomposition précédente (la dimension relative
se comportant bien par changement de base ; voir Morphismes d'espaces, lemme
\ref{spaces-morphisms-lemma-dimension-fibre-after-base-change}).
La méthode exposée à la section \ref{section-open} fournit donc le sous-champ
ouvert possédant les propriétés voulues.
\end{proof}





\section{Courbes de genre donné}
\label{section-genus}

\noindent
Dans le projet Champs, on convient que le genre $g$ d'un
schéma $X$ propre de dimension $1$ sur un corps $k$ n'est défini que
si $H^0(X, \mathcal{O}_X) = k$. Dans ce cas,
$g = \dim_k H^1(X, \mathcal{O}_X)$.
Voir Courbes algébriques, section \ref{curves-section-genus}.
Les conditions nécessaires pour définir le genre définissent un sous-champ ouvert,
qui est alors réunion disjointe de sous-champs ouverts, un pour chaque genre.

\begin{lemma}
\label{lemma-pre-genus-curves}
Il existe un sous-champ ouvert $\Curvesstack^{h0, 1} \subset \Curvesstack$
tel que :
\begin{enumerate}
\item pour une famille de courbes $f : X \to S$, les assertions suivantes
sont équivalentes :
\begin{enumerate}
\item le morphisme classifiant $S \to \Curvesstack$ se factorise par
$\Curvesstack^{h0, 1}$,
\item $f_*\mathcal{O}_X = \mathcal{O}_S$, cette égalité reste vraie
après tout changement de base, et les fibres de $f$ sont de dimension $1$ ;
\end{enumerate}
\item pour un schéma $X$ propre sur un corps $k$ et tel que $\dim(X) \leq 1$,
les assertions suivantes sont équivalentes :
\begin{enumerate}
\item le morphisme classifiant $\Spec(k) \to \Curvesstack$ se factorise par
$\Curvesstack^{h0, 1}$,
\item $H^0(X, \mathcal{O}_X) = k$ et $\dim(X) = 1$.
\end{enumerate}
\end{enumerate}
\end{lemma}

\begin{proof}
Pour une famille de courbes $X \to S$, l'ensemble des $s \in S$ tels que
$\kappa(s) = H^0(X_s, \mathcal{O}_{X_s})$
est ouvert dans $S$ d'après Catégories dérivées des espaces, lemme
\ref{spaces-perfect-lemma-jump-loci-geometric}.
De même, l'ensemble des points de $S$ où la fibre est de
dimension $1$ est ouvert d'après Compléments sur les morphismes d'espaces, lemme
\ref{spaces-more-morphisms-lemma-dimension-fibres-proper-flat}.
En outre, si $f : X \to S$ est une famille de courbes dont toutes les fibres
sont de dimension $1$ (et, en particulier, $f$ est surjectif), alors
la condition (1)(b) équivaut à
$\kappa(s) = H^0(X_s, \mathcal{O}_{X_s})$ pour tout $s \in S$ ; voir
Catégories dérivées des espaces, lemme \ref{spaces-perfect-lemma-proper-flat-h0}.
Le lemme résulte donc de la discussion générale de la
section \ref{section-open}.
\end{proof}

\begin{lemma}
\label{lemma-pre-genus-in-CM-1}
On a $\Curvesstack^{h0, 1} \subset \Curvesstack^{CM, 1}$
en tant que sous-champs ouverts de $\Curvesstack$.
\end{lemma}

\begin{proof}
Voir Courbes algébriques, lemme \ref{curves-lemma-automatic}, ainsi que les
lemmes \ref{lemma-pre-genus-curves} et \ref{lemma-CM-1-curves}.
\end{proof}

\begin{lemma}
\label{lemma-genus}
Soit $f : X \to S$ une famille de courbes telle que
$\kappa(s) = H^0(X_s, \mathcal{O}_{X_s})$ pour tout $s \in S$, c'est-à-dire que
le morphisme classifiant $S \to \Curvesstack$ se factorise par
$\Curvesstack^{h0, 1}$ (lemme \ref{lemma-pre-genus-curves}). Alors :
\begin{enumerate}
\item $f_*\mathcal{O}_X = \mathcal{O}_S$, et cette égalité subsiste après tout changement de base ;
\item $R^1f_*\mathcal{O}_X$ est un $\mathcal{O}_S$-module fini localement libre ;
\item pour tout morphisme $h : S' \to S$, si $f' : X' \to S'$ est obtenu par
changement de base, alors $h^*(R^1f_*\mathcal{O}_X) = R^1f'_*\mathcal{O}_{X'}$.
\end{enumerate}
\end{lemma}

\begin{proof}
Appliquons Catégories dérivées des espaces, lemme
\ref{spaces-perfect-lemma-proper-flat-h0}.
Cela prouve l'assertion (1). Cela implique aussi que, localement sur $S$,
on peut écrire $Rf_*\mathcal{O}_X = \mathcal{O}_S \oplus P$,
où $P$ est parfait et de Tor-amplitude contenue dans $[1, \infty)$.
Rappelons que la formation de $Rf_*\mathcal{O}_X$ commute
à tout changement de base
(Catégories dérivées des espaces, lemme
\ref{spaces-perfect-lemma-flat-proper-perfect-direct-image-general}).
Ainsi, pour $s \in S$, on a
$$
H^i(P \otimes_{\mathcal{O}_S}^\mathbf{L} \kappa(s)) =
H^i(X_s, \mathcal{O}_{X_s})
\text{ pour }i \geq 1
$$
Ce groupe est nul sauf si $i = 1$, car $X_s$ est un schéma noethérien
de dimension $1$ ; voir
Cohomologie, proposition \ref{cohomology-proposition-vanishing-Noetherian}.
On a alors $P = H^1(P)[-1]$, et $H^1(P)$ est fini localement libre,
par exemple d'après Compléments sur l'algèbre, lemme
\ref{more-algebra-lemma-lift-perfect-from-residue-field}.
Comme tout est compatible avec le changement de base, on voit
également que l'assertion (3) est vraie.
\end{proof}

\begin{lemma}
\label{lemma-pre-genus-one-piece-per-genus}
Il existe une décomposition en sous-champs ouverts et fermés
$$
\Curvesstack^{h0, 1} = \coprod\nolimits_{g \geq 0} \Curvesstack_g
$$
où chaque $\Curvesstack_g$ se caractérise comme suit :
\begin{enumerate}
\item pour une famille de courbes $f : X \to S$, les assertions suivantes
sont équivalentes :
\begin{enumerate}
\item le morphisme classifiant $S \to \Curvesstack$ se factorise par
$\Curvesstack_g$,
\item $f_*\mathcal{O}_X = \mathcal{O}_S$, cette égalité reste vraie après
tout changement de base, les fibres de $f$ sont de dimension $1$, et
$R^1f_*\mathcal{O}_X$ est un $\mathcal{O}_S$-module localement libre de rang $g$ ;
\end{enumerate}
\item pour un schéma $X$ propre sur un corps $k$ et tel que $\dim(X) \leq 1$,
les assertions suivantes sont équivalentes :
\begin{enumerate}
\item le morphisme classifiant $\Spec(k) \to \Curvesstack$ se factorise par
$\Curvesstack_g$,
\item $\dim(X) = 1$, $k = H^0(X, \mathcal{O}_X)$ et
le genre de $X$ est $g$.
\end{enumerate}
\end{enumerate}
\end{lemma}

\begin{proof}
Nous disposons déjà de $\Curvesstack^{h0, 1}$ comme sous-champ ouvert de
$\Curvesstack$, caractérisé par les conditions du lemme qui ne font intervenir
ni $R^1f_*$ ni $H^1$ ; voir le lemme \ref{lemma-pre-genus-curves}.
L'existence de la décomposition en sous-champs ouverts et fermés
résulte immédiatement de la discussion de la section \ref{section-open}
et du lemme \ref{lemma-genus}. Cela prouve la caractérisation de (1).
La caractérisation de (2) résulte de la définition du
genre donnée dans Courbes algébriques, définition \ref{curves-definition-genus}.
\end{proof}









\section{Courbes géométriquement réduites}
\label{section-geometrically-reduced}

\noindent
Il existe un sous-champ ouvert de $\Curvesstack$ qui paramètre
les « courbes » géométriquement réduites.

\begin{lemma}
\label{lemma-geometrically-reduced-curves}
Il existe un sous-champ ouvert $\Curvesstack^{geomred} \subset \Curvesstack$
tel que :
\begin{enumerate}
\item pour une famille de courbes $X \to S$, les assertions suivantes
sont équivalentes :
\begin{enumerate}
\item le morphisme classifiant $S \to \Curvesstack$ se factorise par
$\Curvesstack^{geomred}$,
\item les fibres du morphisme $X \to S$ sont géométriquement réduites
(Compléments sur les morphismes d'espaces, définition
\ref{spaces-more-morphisms-definition-geometrically-reduced-fibre}) ;
\end{enumerate}
\item pour un schéma $X$ propre sur un corps $k$ et tel que $\dim(X) \leq 1$,
les assertions suivantes sont équivalentes :
\begin{enumerate}
\item le morphisme classifiant $\Spec(k) \to \Curvesstack$ se factorise par
$\Curvesstack^{geomred}$,
\item $X$ est géométriquement réduit sur $k$.
\end{enumerate}
\end{enumerate}
\end{lemma}

\begin{proof}
Soit $f : X \to S$ une famille de courbes. D'après
Compléments sur les morphismes d'espaces, lemme
\ref{spaces-more-morphisms-lemma-geometrically-reduced-open},
l'ensemble
$$
E = \{s \in S : \text{la fibre de }X \to S\text{ en }s
\text{ est géométriquement réduite}\}
$$
est ouvert dans $S$. La formation de cet ouvert commute à tout
changement de base d'après
Compléments sur les morphismes d'espaces, lemme
\ref{spaces-more-morphisms-lemma-base-change-fibres-geometrically-reduced}.
La méthode exposée à la section \ref{section-open} fournit donc le sous-champ
ouvert possédant les propriétés voulues.
\end{proof}

\begin{lemma}
\label{lemma-geomred-in-CM}
On a $\Curvesstack^{geomred} \subset \Curvesstack^{CM}$
en tant que sous-champs ouverts de $\Curvesstack$.
\end{lemma}

\begin{proof}
Cela tient à ce qu'un schéma noethérien réduit de
dimension $\leq 1$ est Cohen-Macaulay. Voir
Algèbre, lemme \ref{algebra-lemma-criterion-reduced}.
\end{proof}






\section{Courbes géométriquement réduites et connexes}
\label{section-geometrically-reduced-connected}

\noindent
Il existe un sous-champ ouvert de $\Curvesstack$ qui paramètre
les « courbes » géométriquement réduites et connexes.
Nous éliminons d'emblée les objets de dimension $0$.

\begin{lemma}
\label{lemma-geometrically-reduced-connected-1-curves}
Il existe un sous-champ ouvert $\Curvesstack^{grc, 1} \subset \Curvesstack$
tel que :
\begin{enumerate}
\item pour une famille de courbes $X \to S$, les assertions suivantes
sont équivalentes :
\begin{enumerate}
\item le morphisme classifiant $S \to \Curvesstack$ se factorise par
$\Curvesstack^{grc, 1}$,
\item les fibres géométriques du morphisme $X \to S$ sont
réduites, connexes et de dimension $1$ ;
\end{enumerate}
\item pour un schéma $X$ propre sur un corps $k$ et tel que $\dim(X) \leq 1$,
les assertions suivantes sont équivalentes :
\begin{enumerate}
\item le morphisme classifiant $\Spec(k) \to \Curvesstack$ se factorise par
$\Curvesstack^{grc, 1}$,
\item $X$ est géométriquement réduit, géométriquement connexe
et de dimension $1$.
\end{enumerate}
\end{enumerate}
\end{lemma}

\begin{proof}
D'après les lemmes \ref{lemma-geometrically-reduced-curves},
\ref{lemma-geomred-in-CM}, \ref{lemma-CM-curves} et \ref{lemma-CM-1-curves},
il est clair que l'on a
$$
\Curvesstack^{grc, 1}
\subset
\Curvesstack^{geomred} \cap \Curvesstack^{CM, 1}
$$
si ce sous-champ existe. Soit $f : X \to S$ une famille de courbes telle que $f$
soit Cohen-Macaulay, ait des fibres géométriquement réduites et
soit de dimension relative $1$. D'après
Compléments sur les morphismes d'espaces, lemme
\ref{spaces-more-morphisms-lemma-stein-factorization-etale},
dans la factorisation de Stein
$$
X \to T \to S
$$
le morphisme $T \to S$ est \'etale. Il en résulte qu'il existe
un sous-schéma ouvert et fermé $S' \subset S$ tel que
$X \times_S S' \to S'$ ait des fibres géométriquement
connexes (dans la décomposition de
Morphismes, lemme \ref{morphisms-lemma-finite-locally-free},
appliquée au morphisme fini localement libre $T \to S$,
cela correspond à $S_1$).
La formation de cet ouvert commute à tout changement de base,
car le nombre de composantes connexes des fibres géométriques
est invariant par changement de base (il est également vrai
que la factorisation de Stein commute au changement de base
dans notre cas particulier, mais nous n'en avons pas besoin pour conclure).
La méthode exposée à la section \ref{section-open} fournit donc le sous-champ
ouvert possédant les propriétés voulues.
\end{proof}

\begin{lemma}
\label{lemma-geomredcon-in-h0-1}
On a $\Curvesstack^{grc, 1} \subset \Curvesstack^{h0, 1}$
en tant que sous-champs ouverts de $\Curvesstack$. En particulier, pour toute
famille de courbes $f : X \to S$
dont les fibres géométriques sont réduites, connexes et de dimension $1$,
$R^1f_*\mathcal{O}_X$ est un $\mathcal{O}_S$-module fini localement libre
dont la formation commute à tout changement de base.
\end{lemma}

\begin{proof}
Cela résulte de Variétés, lemme
\ref{varieties-lemma-proper-geometrically-reduced-global-sections},
ainsi que des lemmes \ref{lemma-pre-genus-curves} et
\ref{lemma-geometrically-reduced-connected-1-curves}.
La dernière assertion résulte du lemme \ref{lemma-genus}.
\end{proof}

\begin{lemma}
\label{lemma-one-piece-per-genus}
Il existe une décomposition en sous-champs ouverts et fermés
$$
\Curvesstack^{grc, 1} = \coprod\nolimits_{g \geq 0} \Curvesstack^{grc, 1}_g
$$
où chaque $\Curvesstack^{grc, 1}_g$ se caractérise comme suit :
\begin{enumerate}
\item pour une famille de courbes $f : X \to S$, les assertions suivantes
sont équivalentes :
\begin{enumerate}
\item le morphisme classifiant $S \to \Curvesstack$ se factorise par
$\Curvesstack^{grc, 1}_g$,
\item les fibres géométriques du morphisme $f : X \to S$ sont
réduites, connexes, de dimension $1$, et
$R^1f_*\mathcal{O}_X$ est un $\mathcal{O}_S$-module localement libre
de rang $g$ ;
\end{enumerate}
\item pour un schéma $X$ propre sur un corps $k$ et tel que $\dim(X) \leq 1$,
les assertions suivantes sont équivalentes :
\begin{enumerate}
\item le morphisme classifiant $\Spec(k) \to \Curvesstack$ se factorise par
$\Curvesstack^{grc, 1}_g$,
\item $X$ est géométriquement réduit, géométriquement connexe,
de dimension $1$ et de genre $g$.
\end{enumerate}
\end{enumerate}
\end{lemma}

\begin{proof}
Première démonstration : posons
$\Curvesstack^{grc, 1}_g = \Curvesstack^{grc, 1} \cap \Curvesstack_g$
et combinons les lemmes \ref{lemma-geomredcon-in-h0-1} et
\ref{lemma-pre-genus-one-piece-per-genus}.
Seconde démonstration :
L'existence de la décomposition en sous-champs ouverts et fermés
résulte immédiatement de la discussion de la section \ref{section-open}
et du lemme \ref{lemma-geomredcon-in-h0-1}.
Cela prouve la caractérisation de (1).
La caractérisation de (2) en résulte également, car
le genre d'un schéma $X/k$ propre de dimension $1$,
géométriquement réduit et connexe, est défini
(Courbes algébriques, définition \ref{curves-definition-genus}, et
Variétés, lemme
\ref{varieties-lemma-proper-geometrically-reduced-global-sections})
et vaut $\dim_k H^1(X, \mathcal{O}_X)$.
\end{proof}


 
 
 
 
 
\section{Courbes de Gorenstein}
\label{section-gorenstein}

\noindent
Il existe un sous-champ ouvert de $\Curvesstack$ qui paramètre
les « courbes » de Gorenstein.

\begin{lemma}
\label{lemma-gorenstein-curves}
Il existe un sous-champ ouvert $\Curvesstack^{Gorenstein} \subset \Curvesstack$
tel que :
\begin{enumerate}
\item pour une famille de courbes $X \to S$, les assertions suivantes sont équivalentes :
\begin{enumerate}
\item le morphisme classifiant $S \to \Curvesstack$ se factorise
par $\Curvesstack^{Gorenstein}$,
\item le morphisme $X \to S$ est de Gorenstein ;
\end{enumerate}
\item pour un schéma $X$ propre sur un corps $k$ et tel que $\dim(X) \leq 1$,
les assertions suivantes sont équivalentes :
\begin{enumerate}
\item le morphisme classifiant $\Spec(k) \to \Curvesstack$ se factorise
par $\Curvesstack^{Gorenstein}$,
\item $X$ est de Gorenstein.
\end{enumerate}
\end{enumerate}
\end{lemma}

\begin{proof}
Soit $f : X \to S$ une famille de courbes. D'après
Compléments sur les morphismes d'espaces, lemme
\ref{spaces-more-morphisms-lemma-flat-finite-presentation-gorenstein-open},
l'ensemble
$$
W = \{x \in |X| : f \text{ est de Gorenstein au point }x\}
$$
est ouvert dans $|X|$, et la formation de cet ouvert commute à tout
changement de base. Puisque $f$ est propre, le sous-ensemble
$$
S' = S \setminus f(|X| \setminus W)
$$
de $S$ est ouvert, et $X \times_S S' \to S'$ est de Gorenstein.
De plus, la formation de $S'$ commute à tout changement de base,
car il en est ainsi de $W$.
La méthode exposée à la section \ref{section-open} fournit donc le
sous-champ ouvert possédant les propriétés voulues.
\end{proof}

\begin{lemma}
\label{lemma-gorenstein-1-curves}
Il existe un sous-champ ouvert
$\Curvesstack^{Gorenstein, 1} \subset \Curvesstack$ tel que :
\begin{enumerate}
\item pour une famille de courbes $X \to S$, les assertions suivantes sont équivalentes :
\begin{enumerate}
\item le morphisme classifiant $S \to \Curvesstack$ se factorise
par $\Curvesstack^{Gorenstein, 1}$,
\item le morphisme $X \to S$ est de Gorenstein et de
dimension relative $1$ (Morphismes d'espaces, définition
\ref{spaces-morphisms-definition-relative-dimension}) ;
\end{enumerate}
\item pour un schéma $X$ propre sur un corps $k$ et tel que $\dim(X) \leq 1$,
les assertions suivantes sont équivalentes :
\begin{enumerate}
\item le morphisme classifiant $\Spec(k) \to \Curvesstack$ se factorise
par $\Curvesstack^{Gorenstein, 1}$,
\item $X$ est de Gorenstein et $X$ est équidimensionnel de
dimension $1$.
\end{enumerate}
\end{enumerate}
\end{lemma}

\begin{proof}
Rappelons qu'un schéma de Gorenstein est Cohen-Macaulay
(Dualité pour les schémas, lemme \ref{duality-lemma-gorenstein-CM})
et rappelons
qu'un morphisme de Gorenstein est un morphisme de Cohen-Macaulay
(Dualité pour les schémas, lemme \ref{duality-lemma-gorenstein-CM-morphism}).
Nous pouvons donc prendre
$\Curvesstack^{Gorenstein, 1}$ égale à l'intersection
de $\Curvesstack^{Gorenstein}$ et de $\Curvesstack^{CM, 1}$
dans $\Curvesstack$, et appliquer les
lemmes \ref{lemma-gorenstein-curves} et \ref{lemma-CM-1-curves}.
\end{proof}






\section{Courbes d'intersection complète locale}
\label{section-lci}

\noindent
Il existe un sous-champ ouvert de $\Curvesstack$ qui paramètre
les « courbes » qui sont des intersections complètes locales.

\begin{lemma}
\label{lemma-lci-curves}
Il existe un sous-champ ouvert $\Curvesstack^{lci} \subset \Curvesstack$
tel que :
\begin{enumerate}
\item pour une famille de courbes $X \to S$, les assertions suivantes sont équivalentes :
\begin{enumerate}
\item le morphisme classifiant $S \to \Curvesstack$ se factorise par
$\Curvesstack^{lci}$,
\item $X \to S$ est un morphisme d'intersection complète locale, et
\item $X \to S$ est un morphisme syntomique.
\end{enumerate}
\item pour un schéma propre $X$ sur un corps $k$, de dimension $\leq 1$,
les assertions suivantes sont équivalentes :
\begin{enumerate}
\item le morphisme classifiant $\Spec(k) \to \Curvesstack$ se factorise
par $\Curvesstack^{lci}$,
\item $X$ est une intersection complète locale sur $k$.
\end{enumerate}
\end{enumerate}
\end{lemma}

\begin{proof}
Rappelons qu'être un morphisme syntomique revient à être plat et à être
un morphisme d'intersection complète locale ; voir
Compléments sur les morphismes d'espaces, lemme \ref{spaces-more-morphisms-lemma-flat-lci}.
Ainsi, (1)(b) est équivalente à (1)(c).
À la section \ref{section-open}, nous avons vu
qu'il suffit de montrer que, pour toute famille de courbes
$f : X \to S$, il existe un sous-schéma ouvert $S' \subset S$
tel que $S' \times_S X \to S'$ soit un morphisme d'intersection complète
locale et que la formation de $S'$ commute à tout changement de base.
Cela découle du résultat plus général
Compléments sur les morphismes d'espaces, lemme \ref{spaces-more-morphisms-lemma-where-lci}.
\end{proof}




\section{Courbes à singularités isolées}
\label{section-curves-isolated}

\noindent
Nous pouvons considérer le sous-champ ouvert de $\Curvesstack$
qui paramètre les « courbes » n'ayant qu'un nombre fini de points
singuliers (ceux-ci peuvent correspondre à des composantes de dimension $0$
dans notre cadre).

\begin{lemma}
\label{lemma-isolated-sings-curves}
Il existe un sous-champ ouvert
$\Curvesstack^{+} \subset \Curvesstack$
tel que :
\begin{enumerate}
\item pour une famille de courbes $X \to S$, les assertions suivantes sont équivalentes :
\begin{enumerate}
\item le morphisme classifiant $S \to \Curvesstack$ se factorise par
$\Curvesstack^{+}$,
\item le lieu singulier de $X \to S$, muni
d'une quelconque/certaine structure de sous-espace fermé, est fini sur $S$.
\end{enumerate}
\item pour un schéma propre $X$ sur un corps $k$, de dimension $\leq 1$,
les assertions suivantes sont équivalentes :
\begin{enumerate}
\item le morphisme classifiant $\Spec(k) \to \Curvesstack$ se factorise
par $\Curvesstack^{+}$,
\item $X \to \Spec(k)$ est lisse sauf en un nombre fini de points.
\end{enumerate}
\end{enumerate}
\end{lemma}

\begin{proof}
Pour démontrer le lemme, il suffit de montrer que, pour toute famille de courbes
$f : X \to S$, il existe un sous-schéma ouvert $S' \subset S$
tel que la fibre de $S' \times_S X \to S'$ possède la propriété (2).
(La formation de l'ouvert commutera alors automatiquement au changement de base.)
Par définition, le lieu $T \subset |X|$ des points où $X \to S$
n'est pas lisse est fermé. Soit $Z \subset X$ le sous-espace fermé
défini en munissant $T$ de la structure réduite induite d'espace algébrique
(Propriétés des espaces, définition
\ref{spaces-properties-definition-reduced-induced-space}).
Si maintenant $s \in S$ est un point tel que $Z_s$ soit fini, il existe
un voisinage ouvert $U_s \subset S$ de $s$ tel que
$Z \cap f^{-1}(U_s) \to U_s$ soit fini ; voir
Compléments sur les morphismes d'espaces, lemme
\ref{spaces-more-morphisms-lemma-proper-finite-fibre-finite-in-neighbourhood}.
Cela démontre le lemme.
\end{proof}




\section{Le lieu lisse du champ des courbes}
\label{section-smooth}

\noindent
Le morphisme
$$
\Curvesstack \longrightarrow \Spec(\mathbf{Z})
$$
est lisse au-dessus d'un sous-champ ouvert maximal ; voir
Morphismes de champs, lemme \ref{stacks-morphisms-lemma-where-smooth}.
Nous voulons donner un critère pour qu'une courbe appartienne à ce lieu.
Pour cela, nous utiliserons un peu de théorie des déformations.

\medskip\noindent
Soit $k$ un corps. Soit $X$ un schéma propre de dimension $\leq 1$ sur $k$.
Choisissons un anneau de Cohen $\Lambda$ pour $k$ ; voir
Algèbre, lemme \ref{algebra-lemma-cohen-rings-exist}.
Nous sommes alors dans la situation décrite dans
Problèmes de déformation, exemple \ref{examples-defos-example-schemes} et
lemme \ref{examples-defos-lemma-schemes-RS}.
Nous obtenons ainsi une catégorie de déformations $\Deformationcategory_X$
sur la catégorie $\mathcal{C}_\Lambda$ des $\Lambda$-algèbres artiniennes
locales de corps résiduel $k$.

\begin{lemma}
\label{lemma-in-smooth-locus}
Dans la situation ci-dessus, les assertions suivantes sont équivalentes :
\begin{enumerate}
\item le morphisme classifiant $\Spec(k) \to \Curvesstack$ se factorise
par l'ouvert où $\Curvesstack \to \Spec(\mathbf{Z})$ est lisse,
\item la catégorie de déformations $\Deformationcategory_X$ est sans obstruction.
\end{enumerate}
\end{lemma}

\begin{proof}
Puisque $\Curvesstack \longrightarrow \Spec(\mathbf{Z})$ est localement
de présentation finie (lemme \ref{lemma-curves-qs-lfp}),
la formation du sous-champ ouvert où
$\Curvesstack \longrightarrow \Spec(\mathbf{Z})$ est lisse commute au changement de base
plat
(Morphismes de champs, lemme \ref{stacks-morphisms-lemma-where-smooth}).
Puisque l'anneau de Cohen $\Lambda$ est plat sur $\mathbf{Z}$,
nous pouvons travailler sur $\Lambda$. Autrement dit, il s'agit de montrer que
$$
\Lambda\text{-}\Curvesstack \longrightarrow \Spec(\Lambda)
$$
est lisse dans un voisinage ouvert du point
$x_0 : \Spec(k) \to \Lambda\text{-}\Curvesstack$
défini par $X/k$ si et seulement si $\Deformationcategory_X$ est sans obstruction.

\medskip\noindent
Le lemme résulte maintenant de
Géométrie des champs, lemme \ref{stacks-geometry-lemma-characterize-smoothness},
et de l'égalité
$$
\Deformationcategory_X =
\mathcal{F}_{\Lambda\text{-}\Curvesstack, k, x_0}
$$
Cette égalité n'est pas tout à fait triviale à établir. En effet, dans le
membre de gauche, la catégorie de déformations classifie toutes les déformations plates
$Y \to \Spec(A)$ de $X$ comme schéma sur $A \in \Ob(\mathcal{C}_\Lambda)$.
Dans le membre de droite, la catégorie de déformations classifie tous les
morphismes plats $Y \to \Spec(A)$ de fibre spéciale $X$,
où $Y$ est un espace algébrique et où
$Y \to \Spec(A)$ est propre, de présentation finie et de
dimension relative $\leq 1$. Puisque $A$ est artinien, nous voyons
que $Y$ est un schéma, par exemple d'après Espaces sur les corps, lemme
\ref{spaces-over-fields-lemma-codim-1-point-in-schematic-locus}.
Il reste donc à montrer ceci : une déformation plate $Y \to \Spec(A)$ de
$X$ comme schéma sur un anneau local artinien $A$ de corps résiduel $k$
est propre, de présentation finie et de dimension relative $\leq 1$.
La dimension relative se définit en termes des fibres et cette propriété vaut donc
automatiquement pour $Y/A$ puisqu'elle vaut pour $X/k$.
Le morphisme $Y \to \Spec(A)$ est propre et localement de présentation finie,
car c'est vrai pour $X \to \Spec(k)$ ; voir
Compléments sur les morphismes, lemme \ref{more-morphisms-lemma-deform-property}.
\end{proof}

\noindent
Voici un « grand » ouvert du champ des courbes qui est contenu
dans le lieu lisse.

\begin{lemma}
\label{lemma-big-smooth-part-curves}
Le sous-champ ouvert
$$
\Curvesstack^{lci+} =
\Curvesstack^{lci} \cap \Curvesstack^{+}
\subset \Curvesstack
$$
possède les propriétés suivantes :
\begin{enumerate}
\item $\Curvesstack^{lci+} \to \Spec(\mathbf{Z})$ est lisse,
\item pour une famille de courbes $X \to S$, les assertions suivantes sont équivalentes :
\begin{enumerate}
\item le morphisme classifiant $S \to \Curvesstack$ se factorise par
$\Curvesstack^{lci+}$,
\item $X \to S$ est un morphisme d'intersection complète locale et
le lieu singulier de $X \to S$, muni d'une quelconque/certaine structure
de sous-espace fermé, est fini sur $S$ ;
\end{enumerate}
\item pour un schéma propre $X$ sur un corps $k$, de dimension $\leq 1$,
les assertions suivantes sont équivalentes :
\begin{enumerate}
\item le morphisme classifiant $\Spec(k) \to \Curvesstack$ se factorise
par $\Curvesstack^{lci+}$,
\item $X$ est une intersection complète locale sur $k$ et
$X \to \Spec(k)$ est lisse sauf en un nombre fini de points.
\end{enumerate}
\end{enumerate}
\end{lemma}

\begin{proof}
Si nous pouvons montrer qu'il existe un sous-champ ouvert $\Curvesstack^{lci+}$
dont les points sont caractérisés par (2), nous voyons alors que
(1) résulte de la combinaison du lemme \ref{lemma-in-smooth-locus} avec
Problèmes de déformation, lemme \ref{examples-defos-lemma-curve-isolated-lci}.
Puisque
$$
\Curvesstack^{lci+} = \Curvesstack^{lci} \cap \Curvesstack^{+}
$$
dans $\Curvesstack$, on conclut par les
lemmes \ref{lemma-lci-curves} et \ref{lemma-isolated-sings-curves}.
\end{proof}




\section{Courbes lisses}
\label{section-smooth-curves}

\noindent
Dans cette section, nous étudions les sous-champs ouverts de $\Curvesstack$
qui paramètrent les ``courbes'' lisses.

\begin{lemma}
\label{lemma-smooth-curves}
Il existe des sous-champs ouverts
$$
\Curvesstack^{smooth, 1} \subset \Curvesstack^{smooth} \subset \Curvesstack
$$
tels que
\begin{enumerate}
\item étant donnée une famille de courbes $f : X \to S$, les conditions suivantes sont équivalentes :
\begin{enumerate}
\item le morphisme classifiant $S \to \Curvesstack$ se factorise
par $\Curvesstack^{smooth}$, respectivement par $\Curvesstack^{smooth, 1}$,
\item $f$ est lisse, respectivement lisse de dimension relative $1$,
\end{enumerate}
\item étant donné un schéma $X$ propre sur un corps $k$ avec
$\dim(X) \leq 1$, les conditions suivantes sont équivalentes :
\begin{enumerate}
\item le morphisme classifiant $\Spec(k) \to \Curvesstack$ se factorise
par $\Curvesstack^{smooth}$, respectivement par $\Curvesstack^{smooth, 1}$,
\item $X$ est lisse sur $k$, respectivement $X$ est lisse sur $k$ et
$X$ est équidimensionnel de dimension $1$.
\end{enumerate}
\end{enumerate}
\end{lemma}

\begin{proof}
Pour démontrer les assertions concernant $\Curvesstack^{smooth}$,
il suffit de montrer qu'étant donnée une famille de courbes
$f : X \to S$, il existe un sous-schéma ouvert $S' \subset S$
tel que $S' \times_S X \to S'$ soit lisse et que la
formation de cet ouvert commute au changement de base.
Nous savons qu'il existe un plus grand ouvert $U \subset X$ tel
que $U \to S$ soit lisse et que la formation de $U$ commute
à tout changement de base, voir
Morphismes d'espaces, lemme \ref{spaces-morphisms-lemma-where-smooth}.
Si $T = |X| \setminus |U|$, alors $f(T)$ est fermé dans $S$ puisque $f$ est propre.
En posant $S' = S \setminus f(T)$, on obtient l'ouvert voulu.

\medskip\noindent
Soit $f : X \to S$ une famille de courbes telle que $f$ soit lisse.
Alors les fibres $X_s$ sont lisses sur $\kappa(s)$ et sont donc
de Cohen-Macaulay (on peut le voir, par exemple, en utilisant
Algèbre, lemmes \ref{algebra-lemma-smooth-over-field} et
\ref{algebra-lemma-lci-CM}). Ainsi, on voit que l'on peut poser
$$
\Curvesstack^{smooth, 1} = \Curvesstack^{smooth} \cap
\Curvesstack^{CM, 1}
$$
et les équivalences voulues résultent de ce qui a déjà été
montré pour $\Curvesstack^{smooth}$ et du lemme \ref{lemma-CM-1-curves}.
\end{proof}

\begin{lemma}
\label{lemma-smooth-curves-smooth}
Le morphisme $\Curvesstack^{smooth} \to \Spec(\mathbf{Z})$ est lisse.
\end{lemma}

\begin{proof}
Cela résulte immédiatement de l'observation que
$\Curvesstack^{smooth} \subset \Curvesstack^{lci+}$
et du lemme \ref{lemma-big-smooth-part-curves}.
\end{proof}

\begin{lemma}
\label{lemma-smooth-curves-h0}
Il existe un sous-champ ouvert
$\Curvesstack^{smooth, h0} \subset \Curvesstack$
tel que
\begin{enumerate}
\item étant donnée une famille de courbes $f : X \to S$, les conditions suivantes sont équivalentes :
\begin{enumerate}
\item le morphisme classifiant $S \to \Curvesstack$ se factorise
par $\Curvesstack^{smooth}$,
\item $f_*\mathcal{O}_X = \mathcal{O}_S$, cette égalité reste vraie après tout changement de base,
et $f$ est lisse de dimension relative $1$,
\end{enumerate}
\item étant donné un schéma $X$ propre sur un corps $k$ avec
$\dim(X) \leq 1$, les conditions suivantes sont équivalentes :
\begin{enumerate}
\item le morphisme classifiant $\Spec(k) \to \Curvesstack$ se factorise
par $\Curvesstack^{smooth, h0}$,
\item $X$ est lisse, $\dim(X) = 1$, et $k = H^0(X, \mathcal{O}_X)$,
\item $X$ est lisse, $\dim(X) = 1$, et $X$ est géométriquement connexe,
\item $X$ est lisse, $\dim(X) = 1$, et $X$ est géométriquement intègre, et
\item $X_{\overline{k}}$ est une courbe lisse.
\end{enumerate}
\end{enumerate}
\end{lemma}

\begin{proof}
Si l'on pose
$$
\Curvesstack^{smooth, h0} = \Curvesstack^{smooth} \cap
\Curvesstack^{h0, 1}
$$
on voit alors que (1) résulte des
lemmes \ref{lemma-pre-genus-curves} et \ref{lemma-smooth-curves}.
En fait, cela donne aussi l'équivalence de (2)(a) et (2)(b).
Pour achever la démonstration, il reste à montrer que
(2)(b) est équivalent à chacune des conditions (2)(c), (2)(d) et (2)(e).

\medskip\noindent
Un schéma lisse sur un corps est géométriquement normal
(Variétés, lemme \ref{varieties-lemma-smooth-geometrically-normal}),
la lissité est préservée par changement de base
(Morphismes, lemme \ref{morphisms-lemma-base-change-smooth}), et
la propriété d'être lisse est locale pour la topologie fpqc sur le but
(Descente, lemme \ref{descent-lemma-descending-property-smooth}).
Compte tenu de cela, l'équivalence de (2)(b), (2)(c), 2(d) et (2)(e)
résulte de Variétés, lemme \ref{varieties-lemma-geometrically-normal-stein}.
\end{proof}

\begin{definition}
\label{definition-deligne-mumford-smooth}
\begin{reference}
\cite{DM}
\end{reference}
Nous notons $\mathcal{M}$ et appelons
{\it champ de modules des courbes propres et lisses}
le champ algébrique
$\Curvesstack^{smooth, h0}$ paramétrant les familles de courbes
introduit au lemme \ref{lemma-smooth-curves-h0}.
Pour $g \geq 0$, nous notons $\mathcal{M}_g$ et appelons
{\it champ de modules des courbes propres et lisses de genre $g$}
le champ algébrique introduit au
lemme \ref{lemma-smooth-one-piece-per-genus}.
\end{definition}

\noindent
Voici l'inévitable lemme.

\begin{lemma}
\label{lemma-smooth-one-piece-per-genus}
Il existe une décomposition en sous-champs ouverts et fermés
$$
\mathcal{M} = \coprod\nolimits_{g \geq 0} \mathcal{M}_g
$$
où chaque $\mathcal{M}_g$ se caractérise comme suit :
\begin{enumerate}
\item étant donnée une famille de courbes $f : X \to S$, les conditions suivantes sont équivalentes :
\begin{enumerate}
\item le morphisme classifiant $S \to \Curvesstack$ se factorise
par $\mathcal{M}_g$,
\item $X \to S$ est lisse, $f_*\mathcal{O}_X = \mathcal{O}_S$,
cette égalité reste vraie après tout changement de base, et $R^1f_*\mathcal{O}_X$
est un $\mathcal{O}_S$-module localement libre de rang $g$,
\end{enumerate}
\item étant donné un schéma $X$ propre sur un corps $k$ avec
$\dim(X) \leq 1$, les conditions suivantes sont équivalentes :
\begin{enumerate}
\item le morphisme classifiant $\Spec(k) \to \Curvesstack$ se factorise
par $\mathcal{M}_g$,
\item $X$ est lisse, $\dim(X) = 1$, $k = H^0(X, \mathcal{O}_X)$,
et $X$ est de genre $g$,
\item $X$ est lisse, $\dim(X) = 1$, $X$ est géométriquement connexe, et
$X$ est de genre $g$,
\item $X$ est lisse, $\dim(X) = 1$, $X$ est géométriquement intègre, et
$X$ est de genre $g$, et
\item $X_{\overline{k}}$ est une courbe lisse de genre $g$.
\end{enumerate}
\end{enumerate}
\end{lemma}

\begin{proof}
Combiner les lemmes \ref{lemma-smooth-curves-h0} et
\ref{lemma-pre-genus-one-piece-per-genus}.
On peut aussi utiliser
le lemme \ref{lemma-one-piece-per-genus}
à la place.
\end{proof}

\begin{lemma}
\label{lemma-smooth-curves-h0-smooth}
Les morphismes $\mathcal{M} \to \Spec(\mathbf{Z})$ et
$\mathcal{M}_g \to \Spec(\mathbf{Z})$
sont lisses.
\end{lemma}

\begin{proof}
Comme $\mathcal{M}$ est un sous-champ ouvert de
$\Curvesstack^{lci+}$, cela résulte du
lemme \ref{lemma-big-smooth-part-curves}.
\end{proof}





\section{Densité des courbes lisses}
\label{section-smooth-is-dense}

\noindent
Le titre de cette section est trompeur, car nous n'affirmons pas que
$\Curvesstack^{smooth}$ est dense dans $\Curvesstack$.
En fait, cela est faux, comme Mumford l'a montré dans \cite{PathologiesIV}.
Cependant, nous verrons que les ``courbes'' lisses sont denses
dans un grand ouvert.

\begin{lemma}
\label{lemma-smooth-dense}
L'inclusion
$$
|\Curvesstack^{smooth}| \subset |\Curvesstack^{lci+}|
$$
est celle d'un sous-ensemble ouvert dense.
\end{lemma}

\begin{proof}
Par la construction même de la topologie sur
$|\Curvesstack^{lci+}|$ dans
Propriétés des champs, section \ref{stacks-properties-section-points},
on constate que $|\Curvesstack^{smooth}|$
est un sous-ensemble ouvert. Soit $\xi \in |\Curvesstack^{lci+}|$ un point.
Il existe alors un corps $k$ et un schéma $X$ sur $k$
tel que $X$ soit propre sur $k$, que $\dim(X) \leq 1$,
que $X$ soit localement d'intersection complète sur $k$, et
que $X$ soit lisse sur $k$ sauf en un nombre fini de points, et
que $\xi$ soit la classe d'équivalence du
morphisme classifiant $\Spec(k) \to \Curvesstack^{lci+}$ déterminé par $X$.
Voir le lemme \ref{lemma-big-smooth-part-curves}.
D'après Problèmes de déformation, lemme
\ref{examples-defos-lemma-smoothing-proper-curve-isolated-lci},
il existe un morphisme plat et projectif $Y \to \Spec(k[[t]])$
dont la fibre générique est lisse et la fibre spéciale est
isomorphe à $X$. Considérons le morphisme classifiant
$$
\Spec(k[[t]]) \longrightarrow \Curvesstack^{lci+}
$$
déterminé par $Y$. L'image du point fermé est $\xi$
et l'image du point générique appartient à $|\Curvesstack^{smooth}|$.
Comme le point générique se spécialise au point fermé dans
$|\Spec(k[[t]])|$, on conclut que $\xi$ appartient à l'adhérence
de $|\Curvesstack^{smooth}|$, comme voulu.
\end{proof}






\section{Courbes nodales}
\label{section-nodal-curves}

\noindent
En géométrie algébrique, les courbes nodales jouent un rôle particulier.
Nous suggérons au lecteur de jeter un rapide coup d'œil à certaines des discussions
de Courbes algébriques, sections \ref{curves-section-nodal} et
\ref{curves-section-families-nodal},
ainsi que de Compléments sur les morphismes d'espaces, section
\ref{spaces-more-morphisms-section-families-nodal}.

\begin{lemma}
\label{lemma-nodal-curves}
Il existe un sous-champ ouvert $\Curvesstack^{nodal} \subset \Curvesstack$
tel que
\begin{enumerate}
\item étant donnée une famille de courbes $f : X \to S$, les conditions suivantes sont équivalentes :
\begin{enumerate}
\item le morphisme classifiant $S \to \Curvesstack$ se factorise
par $\Curvesstack^{nodal}$,
\item $f$ est au plus nodal de dimension relative $1$,
\end{enumerate}
\item étant donné un schéma $X$ propre sur un corps $k$ avec
$\dim(X) \leq 1$, les conditions suivantes sont équivalentes :
\begin{enumerate}
\item le morphisme classifiant $\Spec(k) \to \Curvesstack$ se factorise
par $\Curvesstack^{nodal}$,
\item les singularités de $X$ sont au plus nodales et $X$
est équidimensionnel de dimension $1$.
\end{enumerate}
\end{enumerate}
\end{lemma}

\begin{proof}
En fait, il suffit de montrer qu'étant donnée une famille de courbes
$f : X \to S$, il existe un sous-schéma ouvert $S' \subset S$
tel que $S' \times_S X \to S'$ soit au plus nodal de dimension relative $1$
et que la formation de $S'$ commute à tout changement de base.
D'après Compléments sur les morphismes d'espaces, lemme
\ref{spaces-more-morphisms-lemma-locus-where-nodal},
il existe un plus grand sous-espace ouvert $X' \subset X$ tel
que $f|_{X'} : X' \to S$ soit au plus nodal de dimension relative $1$.
De plus, la formation de $X'$ commute au changement de base.
On peut donc prendre
$$
S' = S \setminus |f|(|X| \setminus |X'|)
$$
Cet ensemble est ouvert, car un morphisme propre est universellement fermé par
définition.
\end{proof}

\begin{lemma}
\label{lemma-nodal-curves-smooth}
Le morphisme $\Curvesstack^{nodal} \to \Spec(\mathbf{Z})$ est lisse.
\end{lemma}

\begin{proof}
Cela résulte immédiatement de l'observation que
$\Curvesstack^{nodal} \subset \Curvesstack^{lci+}$
et du lemme \ref{lemma-big-smooth-part-curves}.
\end{proof}




\section{Le faisceau dualisant relatif}
\label{section-relative-dualizing}

\noindent
Cette section sert principalement
à introduire des notations dans le cas des familles de courbes.
La majeure partie du travail a déjà été faite dans le chapitre sur la dualité.

\medskip\noindent
Soit $f : X \to S$ une famille de courbes. Il existe un objet
$\omega_{X/S}^\bullet$ de $D_\QCoh(\mathcal{O}_X)$,
appelé le {\it complexe dualisant relatif}, qui possède la propriété
suivante : pour tout diagramme de changement de base
$$
\xymatrix{
X_U \ar[d]_{f'} \ar[r]_{g'} & X \ar[d]^f \\
U \ar[r]^g & S
}
$$
où $U = \Spec(A)$ est affine, le complexe
$\omega_{X_U/U}^\bullet = L(g')^*\omega_{X/S}^\bullet$
représente le foncteur
$$
D_\QCoh(\mathcal{O}_{X_U}) \longrightarrow \text{Mod}_A,\quad
K \longmapsto \Hom_U(Rf_*K, \mathcal{O}_U)
$$
Plus précisément, soit $(\omega_{X/S}^\bullet, \tau)$
le complexe dualisant relatif de la famille défini dans
Dualité pour les espaces, définition
\ref{spaces-duality-definition-relative-dualizing-proper-flat}.
Son existence est démontrée dans Dualité pour les espaces, lemme
\ref{spaces-duality-lemma-existence-relative-dualizing}.
De plus, la formation de $(\omega_{X/S}^\bullet, \tau)$ commute
à tout changement de base (essentiellement par définition ; une référence
précise est Dualité pour les espaces, lemme
\ref{spaces-duality-lemma-base-change-relative-dualizing}).
Dorénavant, nous identifierons le changement de base de
$\omega_{X/S}^\bullet$ au complexe dualisant relatif
de la famille obtenue par changement de base, sans plus le préciser.

\medskip\noindent
Soit $\{S_i \to S\}$ un recouvrement étale avec les $S_i$ affines tel que
$X_i = X \times_S S_i$ soit un schéma ; voir le lemme \ref{lemma-etale-locally-scheme}.
D'après Dualité pour les espaces, lemme \ref{spaces-duality-lemma-compare},
on constate que $\omega_{X_i/S_i}^\bullet$ s'identifie
au complexe dualisant relatif du morphisme de schémas propre, plat et
de présentation finie $f_i : X_i \to S_i$ étudié dans
Dualité pour les schémas, remarque
\ref{duality-remark-relative-dualizing-complex}.
Ainsi, pour démontrer une propriété de $\omega_{X/S}^\bullet$
locale pour la topologie étale, nous pouvons supposer que $X \to S$ est un morphisme de schémas
et utiliser la théorie développée dans le chapitre sur la dualité pour les schémas.
Plus généralement, après tout changement de base de $X$ qui est un schéma,
le complexe dualisant relatif s'identifie au
complexe dualisant relatif de Dualité pour les schémas, remarque
\ref{duality-remark-relative-dualizing-complex}.
Dorénavant, nous utiliserons cette identification sans plus le préciser.

\medskip\noindent
En particulier, soit $\Spec(k) \to S$ un morphisme, où $k$ est un corps.
Notons $X_k$ le changement de base (c'est un schéma d'après
Espaces sur les corps, lemme
\ref{spaces-over-fields-lemma-codim-1-point-in-schematic-locus}).
Alors $\omega_{X_k/k}^\bullet$ est isomorphe
au complexe $\omega_{X_k}^\bullet$ de
Courbes algébriques, lemme \ref{curves-lemma-duality-dim-1}
(tous deux représentent le même foncteur, de sorte que l'on peut appliquer le lemme de Yoneda,
mais cela résulte en fait des remarques précédentes).
On en conclut que les faisceaux de cohomologie
$H^i(\omega_{X_k/k}^\bullet)$
ne sont non nuls que pour $i = 0, -1$.
Si $X_k$ est de Cohen-Macaulay et équidimensionnel de dimension $1$,
alors on n'a plus que $H^{-1}$ et, si $X_k$ est en outre de Gorenstein,
$H^{-1}(\omega_{X_k/k})$ est inversible ; voir
Courbes algébriques, lemmes \ref{curves-lemma-duality-dim-1-CM}
et \ref{curves-lemma-rr}.

\begin{lemma}
\label{lemma-CM-dualizing}
Soit $X \to S$ une famille de courbes dont les fibres sont de Cohen-Macaulay
et équidimensionnelles de dimension $1$ (lemme \ref{lemma-CM-1-curves}).
Alors $\omega_{X/S}^\bullet = \omega_{X/S}[1]$, où $\omega_{X/S}$
est un $\mathcal{O}_X$-module pseudo-cohérent et plat sur $S$, dont la
formation commute à tout changement de base.
\end{lemma}

\begin{proof}
Nous invitons vivement le lecteur à déduire cet énoncé directement de l'étude ci-dessus
du comportement après changement de base à un corps. Notre démonstration
emploiera une réduction quelque peu laborieuse au cas des schémas noethériens.

\medskip\noindent
Dès que nous aurons montré $\omega_{X/S}^\bullet = \omega_{X/S}[1]$ avec
$\omega_{X/S}$ plat sur $S$, l'assertion relative au changement de base
en découlera, puisque nous savons déjà que la formation de $\omega_{X/S}^\bullet$
commute à tout changement de base. De plus, la pseudo-cohérence
sera automatique, car $\omega_{X/S}^\bullet$ est pseudo-cohérent
par définition. L'annulation des autres faisceaux de cohomologie et la platitude
peuvent se vérifier localement pour la topologie étale. Nous pouvons donc supposer que $f : X \to S$
est un morphisme de schémas avec $S$ affine (voir la discussion précédente).
Écrivons $S = \lim S_i$ comme limite projective filtrante de schémas affines $S_i$
de type fini sur $\mathbf{Z}$.
Puisque $\Curvesstack^{CM, 1}$ est localement de présentation finie sur
$\mathbf{Z}$ (comme sous-champ ouvert de $\Curvesstack$ ; voir les
lemmes \ref{lemma-CM-1-curves} et \ref{lemma-curves-qs-lfp}),
on peut trouver un $i$ et une famille
de courbes $X_i \to S_i$ dont l'image réciproque est $X \to S$
(Limites de champs, lemme
\ref{stacks-limits-lemma-representable-by-spaces-limit-preserving}).
Quitte à augmenter $i$, on peut supposer que $X_i$ est un schéma ;
voir Limites d'espaces, lemme \ref{spaces-limits-lemma-limit-is-scheme}.
Puisque la formation de $\omega_{X/S}^\bullet$ commute à
tout changement de base, on peut remplacer $S$ par $S_i$.
Après ce remplacement, nous pouvons supposer, et nous le faisons, que $S_i$ est noethérien.
Alors $f$ est manifestement un morphisme de Cohen-Macaulay
(Compléments sur les morphismes, définition \ref{more-morphisms-definition-CM}),
en vertu de notre hypothèse sur les fibres.
De plus, $\omega_{X/S}^\bullet = f^!\mathcal{O}_S$
par la construction même de $f^!$ dans
Dualité pour les schémas, section \ref{duality-section-upper-shriek}.
Le lemme résulte donc de Dualité pour les schémas, lemme
\ref{duality-lemma-affine-flat-Noetherian-CM}.
\end{proof}

\begin{definition}
\label{definition-relative-dualizing-sheaf}
Soit $f : X \to S$ une famille de courbes dont les fibres sont de Cohen-Macaulay
et équidimensionnelles de dimension $1$ (lemme \ref{lemma-CM-1-curves}).
Le $\mathcal{O}_X$-module
$$
\omega_{X/S} = H^{-1}(\omega_{X/S}^\bullet)
$$
étudié dans le lemme \ref{lemma-CM-dualizing}
est appelé le {\it faisceau dualisant relatif} de $f$.
\end{definition}

\noindent
Dans la situation de la définition \ref{definition-relative-dualizing-sheaf},
le faisceau dualisant relatif $\omega_{X/S}$ possède la propriété suivante
(qui le caractérise en outre localement sur $S$) :
pour tout diagramme de changement de base
$$
\xymatrix{
X_U \ar[d]_{f'} \ar[r]_{g'} & X \ar[d]^f \\
U \ar[r]^g & S
}
$$
où $U = \Spec(A)$ est affine, le module $\omega_{X_U/U} = (g')^*\omega_{X/S}$
représente le foncteur
$$
\QCoh(\mathcal{O}_{X_U}) \longrightarrow \text{Mod}_A,\quad
\mathcal{F} \longmapsto \Hom_A(H^1(X, \mathcal{F}), A)
$$
Cela résulte immédiatement de la propriété correspondante du complexe
dualisant relatif donnée ci-dessus. En particulier, si $A = k$ est un corps,
on retrouve le module dualisant de $X_k$ introduit et étudié dans
Courbes algébriques, lemmes \ref{curves-lemma-duality-dim-1},
\ref{curves-lemma-duality-dim-1-CM} et \ref{curves-lemma-rr}.

\begin{lemma}
\label{lemma-gorenstein-dualizing}
Soit $X \to S$ une famille de courbes dont les fibres sont de Gorenstein
et équidimensionnelles de dimension $1$ (lemme \ref{lemma-gorenstein-1-curves}).
Alors le faisceau dualisant relatif $\omega_{X/S}$ est un
$\mathcal{O}_X$-module inversible dont la
formation commute à tout changement de base.
\end{lemma}

\begin{proof}
Cela résulte de ce que l'image réciproque du module dualisant relatif
sur une fibre est inversible, d'après la discussion précédente. On
peut aussi raisonner exactement comme dans la démonstration du
lemme \ref{lemma-CM-dualizing} et déduire le résultat de
Dualité pour les schémas, lemme
\ref{duality-lemma-affine-flat-Noetherian-gorenstein}.
\end{proof}









\section{Courbes préstables}
\label{section-prestable-curves}

\noindent
La définition suivante est équivalente à ce qui semble être la
notion généralement admise de famille préstable de courbes.

\begin{definition}
\label{definition-prestable}
Soit $f : X \to S$ une famille de courbes. On dit que $f$ est une
{\it famille préstable de courbes} si
\begin{enumerate}
\item $f$ est au plus nodal de dimension relative $1$, et
\item $f_*\mathcal{O}_X = \mathcal{O}_S$, et cette égalité subsiste après
tout changement de base\footnote{En fait, il suffit d'exiger
$f_*\mathcal{O}_X = \mathcal{O}_S$, car la factorisation de Stein
de $f$ est étale dans ce cas ; voir
Compléments sur les morphismes d'espaces, lemme
\ref{spaces-more-morphisms-lemma-stein-factorization-etale}.
On peut aussi remplacer cette condition en demandant que les fibres
géométriques soient connexes ; voir le lemme \ref{lemma-geomredcon-in-h0-1}.}.
\end{enumerate}
\end{definition}

\noindent
Soit $X$ un schéma propre sur un corps $k$ avec $\dim(X) \leq 1$.
Alors $X \to \Spec(k)$ est une famille de courbes ; on peut donc demander
si elle est préstable ou non\footnote{Nous ne pouvons pas employer le terme
``courbe préstable'' ici, car le mot courbe implique l'irréductibilité. Voir
la discussion dans Courbes algébriques, section \ref{curves-section-families-nodal}.}
au sens de la définition. En explicitant les définitions, on voit que
les conditions suivantes sont équivalentes :
\begin{enumerate}
\item $X$ est préstable,
\item les singularités de $X$ sont au plus nodales, $\dim(X) = 1$,
et $k = H^0(X, \mathcal{O}_X)$,
\item $X_{\overline{k}}$ est connexe et lisse sur $\overline{k}$
en dehors d'un nombre fini de nœuds
(Courbes algébriques, définition \ref{curves-definition-multicross}).
\end{enumerate}
Cela montre que notre définition concorde avec la plupart de celles que l'on trouve
dans la littérature.

\begin{lemma}
\label{lemma-prestable-curves}
Il existe un sous-champ ouvert $\Curvesstack^{prestable} \subset \Curvesstack$
tel que
\begin{enumerate}
\item étant donnée une famille de courbes $f : X \to S$, les conditions suivantes sont équivalentes :
\begin{enumerate}
\item le morphisme classifiant $S \to \Curvesstack$ se factorise
par $\Curvesstack^{prestable}$,
\item $X \to S$ est une famille préstable de courbes,
\end{enumerate}
\item étant donné un schéma $X$ propre sur un corps $k$ avec
$\dim(X) \leq 1$, les conditions suivantes sont équivalentes :
\begin{enumerate}
\item le morphisme classifiant $\Spec(k) \to \Curvesstack$
se factorise par $\Curvesstack^{prestable}$,
\item les singularités de $X$ sont au plus nodales, $\dim(X) = 1$,
et $k = H^0(X, \mathcal{O}_X)$.
\end{enumerate}
\end{enumerate}
\end{lemma}

\begin{proof}
Étant donnée une famille de courbes $X \to S$, elle est préstable si
et seulement si le morphisme classifiant se factorise à la fois par
$\Curvesstack^{nodal}$ et $\Curvesstack^{h0, 1}$. On peut aussi
utiliser $\Curvesstack^{grc, 1}$ (puisqu'une courbe nodale est géométriquement
réduite et a donc son $H^0$ égal au corps de base si et seulement si
elle est connexe). En formule,
$$
\Curvesstack^{prestable} =
\Curvesstack^{nodal} \cap \Curvesstack^{h0, 1} =
\Curvesstack^{nodal} \cap \Curvesstack^{grc, 1}
$$
Le lemme résulte donc des
lemmes \ref{lemma-pre-genus-curves} et \ref{lemma-nodal-curves}.
\end{proof}

\noindent
Pour tout genre $g \geq 0$, nous disposons du champ algébrique classifiant
les courbes préstables de genre $g$. En fait, dorénavant, nous dirons
que $X \to S$ est une {\it famille préstable de courbes de genre $g$}
si et seulement si le morphisme classifiant $S \to \Curvesstack$ se factorise par
le sous-champ ouvert $\Curvesstack^{prestable}_g$ du
lemme \ref{lemma-prestable-one-piece-per-genus}.

\begin{lemma}
\label{lemma-prestable-one-piece-per-genus}
Il existe une décomposition en sous-champs ouverts et fermés
$$
\Curvesstack^{prestable} = \coprod\nolimits_{g \geq 0}
\Curvesstack^{prestable}_g
$$
où chaque $\Curvesstack^{prestable}_g$ est caractérisé de la manière suivante :
\begin{enumerate}
\item étant donnée une famille de courbes $f : X \to S$, les conditions suivantes sont équivalentes :
\begin{enumerate}
\item le morphisme classifiant $S \to \Curvesstack$ se factorise
par $\Curvesstack^{prestable}_g$,
\item $X \to S$ est une famille préstable de courbes et
$R^1f_*\mathcal{O}_X$ est un $\mathcal{O}_S$-module localement libre de rang $g$,
\end{enumerate}
\item étant donné un schéma $X$ propre sur un corps $k$ avec
$\dim(X) \leq 1$, les conditions suivantes sont équivalentes :
\begin{enumerate}
\item le morphisme classifiant $\Spec(k) \to \Curvesstack$
se factorise par $\Curvesstack^{prestable}_g$,
\item les singularités de $X$ sont au plus nodales, $\dim(X) = 1$,
$k = H^0(X, \mathcal{O}_X)$, et le genre de $X$ est $g$.
\end{enumerate}
\end{enumerate}
\end{lemma}

\begin{proof}
Puisque nous avons vu que $\Curvesstack^{prestable}$ est contenu
dans $\Curvesstack^{h0, 1}$, l'assertion résulte des
lemmes \ref{lemma-prestable-curves} et
\ref{lemma-pre-genus-one-piece-per-genus}.
\end{proof}

\begin{lemma}
\label{lemma-prestable-curves-smooth}
Les morphismes
$\Curvesstack^{prestable} \to \Spec(\mathbf{Z})$ et
$\Curvesstack^{prestable}_g \to \Spec(\mathbf{Z})$ sont
lisses.
\end{lemma}

\begin{proof}
Puisque $\Curvesstack^{prestable}$ est un sous-champ ouvert de
$\Curvesstack^{nodal}$, cela résulte du
lemme \ref{lemma-nodal-curves-smooth}.
\end{proof}




\section{Courbes semi-stables}
\label{section-semistable-curves}

\noindent
Le lemme suivant nous aidera à comprendre les familles de courbes semi-stables.

\begin{lemma}
\label{lemma-semistable}
Soit $f : X \to S$ une famille préstable de courbes de genre $g \geq 1$.
Soit $s \in S$ un point du schéma de base. Soit $m \geq 2$.
Les conditions suivantes sont équivalentes :
\begin{enumerate}
\item $X_s$ n'a pas de queue rationnelle
(Courbes algébriques, exemple \ref{curves-example-rational-tail}), et
\item $f^*f_*\omega_{X/S}^{\otimes m} \to \omega_{X/S}^{\otimes m}$
est surjectif sur $f^{-1}(U)$ pour un ouvert $s \in U \subset S$.
\end{enumerate}
\end{lemma}

\begin{proof}
Supposons (2). D'après les résultats de la section \ref{section-relative-dualizing},
on en déduit que $\omega_{X_s}^{\otimes m}$ est
globalement engendré. Cependant, si $C \subset X_s$
est une queue rationnelle, alors $\deg(\omega_{X_s}|_C) < 0$ d'après
Courbes algébriques, lemme \ref{curves-lemma-rational-tail-negative},
d'où $H^0(C, \omega_{X_s}|_C) = 0$ d'après
Variétés, lemme \ref{varieties-lemma-check-invertible-sheaf-trivial},
ce qui contredit le fait qu'il soit globalement engendré.
Cela prouve (1).

\medskip\noindent
Supposons (1). Supposons d'abord que $g \geq 2$. L'hypothèse (1)
implique que $\omega_{X_s}^{\otimes m}$ est globalement engendré ;
voir Courbes algébriques, lemme \ref{curves-lemma-contracting-rational-tails}.
En outre, on a
$$
\Hom_{\kappa(s)}(H^1(X_s, \omega_{X_s}^{\otimes m}), \kappa(s)) =
H^0(X_s, \omega_{X_s}^{\otimes 1 - m})
$$
par dualité ; voir Courbes algébriques, lemme \ref{curves-lemma-duality-dim-1-CM}.
Puisque $\omega_{X_s}^{\otimes m}$ est globalement engendré, on voit
que sa restriction à chaque composante irréductible est de degré non négatif.
Par conséquent, la restriction de $\omega_{X_s}^{\otimes 1 - m}$ à chaque
composante irréductible est de degré non positif. Comme
$\deg(\omega_{X_s}^{\otimes 1 - m}) = (1 - m)(2g - 2) < 0$ par Riemann--Roch
(Courbes algébriques, lemme \ref{curves-lemma-rr}), on conclut que ce $H^0$
est nul d'après Variétés, lemme \ref{varieties-lemma-no-sections-dual-nef}.
Le théorème de cohomologie et changement de base montre que
$$
E = Rf_*\omega_{X/S}^{\otimes m}
$$
est un complexe parfait dont la formation commute à tout changement de base
(Catégories dérivées des espaces, lemme
\ref{spaces-perfect-lemma-flat-proper-perfect-direct-image-general}).
L'annulation démontrée ci-dessus montre que $E \otimes^\mathbf{L} \kappa(s)$
est égal à $H^0(X_s, \omega_{X_s}^{\otimes m})$ placé en degré $0$.
Quitte à restreindre $S$, on voit que $E = f_*\omega_{X/S}^{\otimes m}$
est un $\mathcal{O}_S$-module localement libre placé en degré $0$
(et que sa formation commute à tout changement de base, comme nous
l'avons déjà dit) ; voir Catégories dérivées des espaces, lemme
\ref{spaces-perfect-lemma-open-where-cohomology-in-degree-i-rank-r-geometric}.
Le morphisme $f^*f_*\omega_{X/S}^{\otimes m} \to \omega_{X/S}^{\otimes m}$
devient surjectif après restriction à $X_s$. Il est donc surjectif dans
un voisinage ouvert de $X_s$. Puisque $f$ est propre, ce voisinage
contient $f^{-1}(U)$ pour un voisinage ouvert
$U$ de $s$ dans $S$.

\medskip\noindent
Supposons (1) et $g = 1$. D'après
Courbes algébriques, lemme \ref{curves-lemma-contracting-rational-tails},
l'hypothèse (1) signifie que $\omega_{X_s}$ est isomorphe à
$\mathcal{O}_{X_s}$. S'il est possible de montrer que, quitte à restreindre $S$,
le faisceau inversible $\omega_{X/S}$ devient trivial, alors
la preuve est achevée. On peut supposer $S$ affine. En restreignant encore $S$,
on peut écrire
$$
Rf_*\mathcal{O}_X = (\mathcal{O}_S \xrightarrow{0} \mathcal{O}_S)
$$
placé en degrés $0$ et $1$,
de manière compatible à tout changement de base ultérieur ; voir le lemme \ref{lemma-genus}.
Par dualité, cela signifie que
$$
Rf_*\omega_{X/S} = (\mathcal{O}_S \xrightarrow{0} \mathcal{O}_S)
$$
placé en degrés $0$ et $1$\footnote{On utilise que
$Rf_*\omega_{X/S}^\bullet =
Rf_*R\SheafHom_{\mathcal{O}_X}(\mathcal{O}_X. \omega_{X/S}^\bullet) =
R\SheafHom_{\mathcal{O}_S}(Rf_*\mathcal{O}_X, \mathcal{O}_S)$
d'après Dualité pour les espaces, lemme \ref{spaces-duality-lemma-iso-on-RSheafHom} et
remarque \ref{spaces-duality-remark-iso-on-RSheafHom},
puis que $\omega_{X/S}^\bullet = \omega_{X/S}[1]$ d'après
nos définitions de la section \ref{section-relative-dualizing}.}.
En particulier, on obtient un isomorphisme $\mathcal{O}_S \to f_*\omega_{X/S}$
qui est compatible à tout changement de base, puisque la
formation de $Rf_*\omega_{X/S}$ est compatible au changement de base
(voir la référence donnée ci-dessus).
Par adjonction, on obtient une section globale $\sigma \in \Gamma(X, \omega_{X/S})$.
La restriction de cette section à la fibre $X_s$
est non nulle (c'est même un élément d'une base) et, puisque
$\omega_{X_s}$ est trivial sur les fibres,
cette section ne s'annule en aucun point de $X_s$.
Elle ne s'annule donc pas dans
un voisinage ouvert de $X_s$. Puisque $f$ est propre, ce voisinage
contient $f^{-1}(U)$ pour un voisinage ouvert
$U$ de $s$ dans $S$.
\end{proof}

\noindent
Le lemme \ref{lemma-semistable} motive la définition suivante.

\begin{definition}
\label{definition-semistable}
Soit $f : X \to S$ une famille de courbes.
On dit que $f$ est une {\it famille semi-stable de courbes} si
\begin{enumerate}
\item $X \to S$ est une famille préstable de courbes, et
\item $X_s$ est de genre $\geq 1$ et
n'a pas de queue rationnelle pour tout $s \in S$.
\end{enumerate}
\end{definition}

\noindent
En particulier, une famille préstable de courbes de genre $0$ n'est jamais
semi-stable.
Soit $X$ un schéma propre sur un corps $k$ avec $\dim(X) \leq 1$.
Alors $X \to \Spec(k)$ est une famille de courbes ; on peut donc demander
si elle est semi-stable ou non. En développant les définitions, on voit que
les conditions suivantes sont équivalentes :
\begin{enumerate}
\item $X$ est semi-stable,
\item $X$ est préstable, est de genre $\geq 1$ et n'a pas de queue rationnelle,
\item $X_{\overline{k}}$ est connexe, est lisse sur $\overline{k}$
en dehors d'un nombre fini de nœuds, est de genre $\geq 1$, et n'a aucune
composante irréductible isomorphe à $\mathbf{P}^1_{\overline{k}}$
qui rencontre le reste de $X_{\overline{k}}$ en un seul point.
\end{enumerate}
Pour voir l'équivalence de (2) et (3), on utilise que $X$ n'a pas de queue rationnelle
si et seulement si $X_{\overline{k}}$ n'a pas de queue rationnelle, d'après
Courbes algébriques, lemme \ref{curves-lemma-contracting-rational-tails}.
Cela montre que notre définition concorde avec la plupart de celles que l'on trouve
dans la littérature.

\begin{lemma}
\label{lemma-semistable-curves}
Il existe un sous-champ ouvert $\Curvesstack^{semistable} \subset \Curvesstack$
tel que
\begin{enumerate}
\item étant donnée une famille de courbes $f : X \to S$, les conditions suivantes sont équivalentes :
\begin{enumerate}
\item le morphisme classifiant $S \to \Curvesstack$ se factorise
par $\Curvesstack^{semistable}$,
\item $X \to S$ est une famille semi-stable de courbes,
\end{enumerate}
\item étant donné un schéma $X$ propre sur un corps $k$ avec
$\dim(X) \leq 1$, les conditions suivantes sont équivalentes :
\begin{enumerate}
\item le morphisme classifiant $\Spec(k) \to \Curvesstack$
se factorise par $\Curvesstack^{semistable}$,
\item les singularités de $X$ sont au plus nodales, $\dim(X) = 1$,
$k = H^0(X, \mathcal{O}_X)$, le genre de $X$ est $\geq 1$, et
$X$ n'a pas de queue rationnelle,
\item les singularités de $X$ sont au plus nodales, $\dim(X) = 1$,
$k = H^0(X, \mathcal{O}_X)$, et $\omega_{X_s}^{\otimes m}$ est
globalement engendré pour $m \geq 2$.
\end{enumerate}
\end{enumerate}
\end{lemma}

\begin{proof}
L'équivalence de (2)(b) et (2)(c) est donnée par
Courbes algébriques, lemme \ref{curves-lemma-contracting-rational-tails}.
Dans la suite de la démonstration, nous travaillerons avec (2)(b),
conformément à la définition \ref{definition-semistable}.

\medskip\noindent
D'après la discussion de la section \ref{section-open},
il suffit de considérer les familles $f : X \to S$ de
courbes préstables. Le lemme \ref{lemma-semistable}
donne l'ouverture souhaitée du lieu en question.
La formation de cet ouvert commute à tout changement de base,
car l'existence ou la non-existence de queues rationnelles est invariante
par extension du corps de base, d'après
Courbes algébriques, lemme \ref{curves-lemma-contracting-rational-tails}.
\end{proof}

\begin{lemma}
\label{lemma-semistable-one-piece-per-genus}
Il existe une décomposition en sous-champs ouverts et fermés
$$
\Curvesstack^{semistable} = \coprod\nolimits_{g \geq 1}
\Curvesstack^{semistable}_g
$$
où chaque $\Curvesstack^{semistable}_g$ est caractérisé de la manière suivante :
\begin{enumerate}
\item étant donnée une famille de courbes $f : X \to S$, les conditions suivantes sont équivalentes :
\begin{enumerate}
\item le morphisme classifiant $S \to \Curvesstack$ se factorise
par $\Curvesstack^{semistable}_g$,
\item $X \to S$ est une famille semi-stable de courbes et
$R^1f_*\mathcal{O}_X$ est un $\mathcal{O}_S$-module localement libre de rang $g$,
\end{enumerate}
\item étant donné un schéma $X$ propre sur un corps $k$ avec
$\dim(X) \leq 1$, les conditions suivantes sont équivalentes :
\begin{enumerate}
\item le morphisme classifiant $\Spec(k) \to \Curvesstack$
se factorise par $\Curvesstack^{semistable}_g$,
\item les singularités de $X$ sont au plus nodales, $\dim(X) = 1$,
$k = H^0(X, \mathcal{O}_X)$, le genre de $X$ est $g$, et $X$
n'a pas de queue rationnelle,
\item les singularités de $X$ sont au plus nodales, $\dim(X) = 1$,
$k = H^0(X, \mathcal{O}_X)$, le genre de $X$ est $g$, et
$\omega_{X_s}^{\otimes m}$ est globalement engendré pour $m \geq 2$.
\end{enumerate}
\end{enumerate}
\end{lemma}

\begin{proof}
Combiner les lemmes \ref{lemma-semistable-curves} et
\ref{lemma-prestable-one-piece-per-genus}.
\end{proof}

\begin{lemma}
\label{lemma-semistable-curves-smooth}
Les morphismes
$\Curvesstack^{semistable} \to \Spec(\mathbf{Z})$ et
$\Curvesstack^{semistable}_g \to \Spec(\mathbf{Z})$
sont lisses.
\end{lemma}

\begin{proof}
Puisque $\Curvesstack^{semistable}$ est un sous-champ ouvert de
$\Curvesstack^{nodal}$, cela résulte du
lemme \ref{lemma-nodal-curves-smooth}.
\end{proof}







\section{Courbes stables}
\label{section-stable-curves}

\noindent
Le lemme suivant nous aidera à comprendre les familles de courbes stables.

\begin{lemma}
\label{lemma-stable}
Soit $f : X \to S$ une famille préstable de courbes de genre $g \geq 2$.
Soit $s \in S$ un point du schéma de base.
Les conditions suivantes sont équivalentes :
\begin{enumerate}
\item $X_s$ n'a ni queue rationnelle ni
pont rationnel
(Courbes algébriques, exemples
\ref{curves-example-rational-tail} et
\ref{curves-example-rational-bridge}) ;
\item $\omega_{X/S}$ est ample sur $f^{-1}(U)$ pour un ouvert $s \in U \subset S$.
\end{enumerate}
\end{lemma}

\begin{proof}
Supposons (2). Alors $\omega_{X_s}$ est ample sur $X_s$.
D'après Courbes algébriques, lemmes \ref{curves-lemma-rational-tail-negative} et
\ref{curves-lemma-rational-bridge-zero},
nous en concluons que (1) est satisfaite (nous utilisons aussi
la caractérisation des faisceaux inversibles amples donnée dans
Variétés, lemme
\ref{varieties-lemma-ampleness-in-terms-of-degrees-components}).

\medskip\noindent
Supposons (1). Alors $\omega_{X_s}$ est ample sur $X_s$ d'après
Courbes algébriques, lemme \ref{curves-lemma-contracting-rational-bridges}.
La conclusion résulte de Descente sur les espaces, lemme
\ref{spaces-descent-lemma-ample-in-neighbourhood}.
\end{proof}

\noindent
Le lemme \ref{lemma-stable} motive la définition suivante.

\begin{definition}
\label{definition-stable}
Soit $f : X \to S$ une famille de courbes.
On dit que $f$ est une {\it famille stable de courbes} si
\begin{enumerate}
\item $X \to S$ est une famille préstable de courbes ;
\item $X_s$ est de genre $\geq 2$ et ne possède ni queue rationnelle
ni pont rationnel pour tout $s \in S$.
\end{enumerate}
\end{definition}

\noindent
En particulier, une famille préstable de courbes de genre $0$ ou $1$ n'est jamais
stable.
Soit $X$ un schéma propre sur un corps $k$ avec $\dim(X) \leq 1$.
Alors $X \to \Spec(k)$ est une famille de courbes ; on peut donc demander
si elle est stable ou non. En explicitant les définitions, on voit que
les conditions suivantes sont équivalentes :
\begin{enumerate}
\item $X$ est stable ;
\item $X$ est préstable, est de genre $\geq 2$ et ne possède ni queue rationnelle
ni pont rationnel ;
\item $X$ est géométriquement connexe, est lisse sur $k$
en dehors d'un nombre fini de nœuds, et $\omega_X$ est ample.
\end{enumerate}
Pour voir l'équivalence de (2) et (3), on utilise
le lemme \ref{lemma-stable} ci-dessus.
Cela montre que notre définition concorde avec la plupart de celles que l'on trouve
dans la littérature.

\begin{lemma}
\label{lemma-stable-curves}
Il existe un sous-champ ouvert $\Curvesstack^{stable} \subset \Curvesstack$
tel que :
\begin{enumerate}
\item pour une famille de courbes $f : X \to S$, les conditions suivantes sont équivalentes :
\begin{enumerate}
\item le morphisme classifiant $S \to \Curvesstack$ se factorise
par $\Curvesstack^{stable}$ ;
\item $X \to S$ est une famille stable de courbes ;
\end{enumerate}
\item pour un schéma $X$ propre sur un corps $k$ avec
$\dim(X) \leq 1$, les conditions suivantes sont équivalentes :
\begin{enumerate}
\item le morphisme classifiant $\Spec(k) \to \Curvesstack$
se factorise par $\Curvesstack^{stable}$ ;
\item les singularités de $X$ sont au plus nodales, $\dim(X) = 1$,
$k = H^0(X, \mathcal{O}_X)$, le genre de $X$ est $\geq 2$, et
$X$ n'a ni queue rationnelle ni pont rationnel ;
\item les singularités de $X$ sont au plus nodales, $\dim(X) = 1$,
$k = H^0(X, \mathcal{O}_X)$, et $\omega_{X_s}$ est ample.
\end{enumerate}
\end{enumerate}
\end{lemma}

\begin{proof}
D'après la discussion de la section \ref{section-open},
il suffit de considérer les familles $f : X \to S$ de
courbes préstables. Le lemme \ref{lemma-stable}
donne l'ouverture souhaitée du lieu considéré.
La formation de cet ouvert commute à tout changement de base,
soit parce que l'existence ou la non-existence de queues ou de ponts rationnels
est invariante par extension du corps de base, d'après
Courbes algébriques, lemmes
\ref{curves-lemma-contracting-rational-tails} et
\ref{curves-lemma-contracting-rational-bridges},
soit parce que l'amplitude est invariante par extension du corps de base, d'après
Descente, lemme \ref{descent-lemma-descending-property-ample}.
\end{proof}

\begin{definition}
\label{definition-deligne-mumford}
\begin{reference}
\cite{DM}
\end{reference}
Nous notons $\overline{\mathcal{M}}$ et appelons
{\it champ de modules des courbes stables} le champ algébrique
$\Curvesstack^{stable}$ qui paramètre les familles stables de courbes
introduit au lemme \ref{lemma-stable-curves}.
Pour $g \geq 2$, nous notons $\overline{\mathcal{M}}_g$ et appelons
{\it champ de modules des courbes stables de genre $g$}
le champ algébrique introduit au lemme \ref{lemma-stable-one-piece-per-genus}.
\end{definition}

\noindent
Voici le lemme incontournable.

\begin{lemma}
\label{lemma-stable-one-piece-per-genus}
Il existe une décomposition en sous-champs ouverts et fermés
$$
\overline{\mathcal{M}} = \coprod\nolimits_{g \geq 2} \overline{\mathcal{M}}_g
$$
où chaque $\overline{\mathcal{M}}_g$ est caractérisé comme suit :
\begin{enumerate}
\item pour une famille de courbes $f : X \to S$, les conditions suivantes sont équivalentes :
\begin{enumerate}
\item le morphisme classifiant $S \to \Curvesstack$ se factorise
par $\overline{\mathcal{M}}_g$ ;
\item $X \to S$ est une famille stable de courbes et
$R^1f_*\mathcal{O}_X$ est un $\mathcal{O}_S$-module localement libre de rang $g$ ;
\end{enumerate}
\item pour un schéma $X$ propre sur un corps $k$ avec
$\dim(X) \leq 1$, les conditions suivantes sont équivalentes :
\begin{enumerate}
\item le morphisme classifiant $\Spec(k) \to \Curvesstack$
se factorise par $\overline{\mathcal{M}}_g$ ;
\item les singularités de $X$ sont au plus nodales, $\dim(X) = 1$,
$k = H^0(X, \mathcal{O}_X)$, le genre de $X$ est $g$, et $X$
n'a ni queue rationnelle ni pont rationnel ;
\item les singularités de $X$ sont au plus nodales, $\dim(X) = 1$,
$k = H^0(X, \mathcal{O}_X)$, le genre de $X$ est $g$, et
$\omega_{X_s}$ est ample.
\end{enumerate}
\end{enumerate}
\end{lemma}

\begin{proof}
Combiner les lemmes \ref{lemma-stable-curves} et
\ref{lemma-prestable-one-piece-per-genus}.
\end{proof}

\begin{lemma}
\label{lemma-stable-curves-smooth}
Les morphismes
$\overline{\mathcal{M}} \to \Spec(\mathbf{Z})$ et
$\overline{\mathcal{M}}_g \to \Spec(\mathbf{Z})$
sont lisses.
\end{lemma}

\begin{proof}
Puisque $\overline{\mathcal{M}}$ est un sous-champ ouvert de
$\Curvesstack^{nodal}$, cela résulte du
lemme \ref{lemma-nodal-curves-smooth}.
\end{proof}

\begin{lemma}
\label{lemma-stable-curves-deligne-mumford}
Les champs $\overline{\mathcal{M}}$ et
$\overline{\mathcal{M}}_g$
sont des sous-champs ouverts de $\Curvesstack^{DM}$.
En particulier, $\overline{\mathcal{M}}$ et
$\overline{\mathcal{M}}_g$ sont DM
(Morphismes de champs, définition
\ref{stacks-morphisms-definition-absolute-separated})
ainsi que des champs de Deligne-Mumford
(Champs algébriques, définition \ref{algebraic-definition-deligne-mumford}).
\end{lemma}

\begin{proof}
Démontrons la première assertion.
Soit $X$ un schéma propre sur un corps $k$ dont les singularités
sont au plus nodales, avec $\dim(X) = 1$, $k = H^0(X, \mathcal{O}_X)$,
tel que le genre de $X$ soit $\geq 2$, et tel que $X$ n'ait
ni queue rationnelle ni pont rationnel.
Il faut montrer que le morphisme classifiant
$\Spec(k) \to \overline{\mathcal{M}} \to \Curvesstack$
se factorise par $\Curvesstack^{DM}$.
Nous pouvons d'abord remplacer $k$ par sa clôture algébrique
(puisque nous savons déjà que les champs considérés sont des sous-champs
ouverts du champ algébrique $\Curvesstack$).
D'après les lemmes \ref{lemma-stable-curves}, \ref{lemma-DM-curves} et
\ref{lemma-in-DM-locus-vector-fields}, il suffit de montrer que
$\text{Der}_k(\mathcal{O}_X, \mathcal{O}_X) = 0$.
Cela est démontré dans
Courbes algébriques, lemme \ref{curves-lemma-stable-vector-fields}.

\medskip\noindent
Puisque $\Curvesstack^{DM}$ est le plus grand sous-champ ouvert de
$\Curvesstack$ qui soit DM, il en va de même du
sous-champ ouvert $\overline{\mathcal{M}}$ de $\Curvesstack^{DM}$.
Enfin, un champ algébrique DM est de Deligne-Mumford d'après
Morphismes de champs, théorème \ref{stacks-morphisms-theorem-DM}.
\end{proof}

\begin{lemma}
\label{lemma-smooth-dense-in-stable}
Soit $g \geq 2$. L'inclusion
$$
|\mathcal{M}_g| \subset |\overline{\mathcal{M}}_g|
$$
est celle d'un sous-ensemble ouvert dense.
\end{lemma}

\begin{proof}
Puisque $\overline{\mathcal{M}}_g \subset \Curvesstack^{lci+}$
est ouvert et que
$\Curvesstack^{smooth} \cap \overline{\mathcal{M}}_g = \mathcal{M}_g$,
cela résulte immédiatement du
lemme \ref{lemma-smooth-dense}.
\end{proof}



\section{Morphismes de contraction}
\label{section-contracting}

\noindent
Nous invitons vivement le lecteur à se familiariser avec
Courbes algébriques, sections
\ref{curves-section-contracting-rational-tails},
\ref{curves-section-contracting-rational-bridges} et
\ref{curves-section-contracting-to-stable},
avant de poursuivre. Le résultat principal de cette section
est l'existence d'un morphisme de « stabilisation »
$$
\Curvesstack^{prestable}_g
\longrightarrow
\overline{\mathcal{M}}_g
$$
Voir le lemme \ref{lemma-stabilization-morphism}.
En termes informels, ce morphisme envoie le point modulaire d'une
courbe nodale de genre $g$ sur le point modulaire de la courbe stable
associée construite dans
Courbes algébriques, lemme \ref{curves-lemma-characterize-contraction-to-stable}.

\begin{lemma}
\label{lemma-contract}
Soient $S$ un schéma et $s \in S$ un point.
Soient $f : X \to S$ et $g : Y \to S$ des familles de courbes.
Soit $c : X \to Y$ un morphisme sur $S$. Si
$c_{s, *}\mathcal{O}_{X_s} = \mathcal{O}_{Y_s}$ et
$R^1c_{s, *}\mathcal{O}_{X_s} = 0$, alors,
après remplacement de $S$ par un voisinage ouvert de $s$,
on a $\mathcal{O}_Y = c_*\mathcal{O}_X$ et $R^1c_*\mathcal{O}_X = 0$,
et cela reste vrai après tout changement de base par un morphisme $S' \to S$.
\end{lemma}

\begin{proof}
Soit $(U, u) \to (S, s)$ un voisinage étale tel que
$\mathcal{O}_{Y_U} = (X_U \to Y_U)_*\mathcal{O}_{X_U}$ et
$R^1(X_U \to Y_U)_*\mathcal{O}_{X_U} = 0$, ces égalités restant vraies
après changement de base par $U' \to U$. Remplaçons alors $S$ par
l'image ouverte de $U \to S$. Étant donné $S' \to S$, posons $U' = U \times_S S'$ ;
nous obtenons les recouvrements étales $\{U' \to S'\}$ et
$\{Y_{U'} \to Y_{S'}\}$. La validité de l'énoncé pour
le changement de base de $c$ par $S' \to S$ résulte donc de sa validité
pour le changement de base de $X_U \to Y_U$ par
$U' \to U$. Autrement dit, la question est locale pour la
topologie étale sur $S$.
Ainsi, d'après le lemme \ref{lemma-etale-locally-scheme}, nous pouvons supposer que
$X$ et $Y$ sont des schémas. D'après
Compléments sur les morphismes, lemme \ref{more-morphisms-lemma-h1-fibre-zero-isom},
il existe un sous-schéma ouvert $V \subset Y$ contenant $Y_s$
tel que $c_*\mathcal{O}_X|_V = \mathcal{O}_V$ et
$R^1c_*\mathcal{O}_X|_V = 0$, ces égalités restant vraies après
tout changement de base par $S' \to S$. Puisque $g : Y \to S$ est propre, on peut
trouver un voisinage ouvert $U \subset S$ de $s$ tel que
$g^{-1}(U) \subset V$. Alors $U$ convient.
\end{proof}

\begin{lemma}
\label{lemma-contract-basic-uniqueness}
Soient $S$ un schéma et $s \in S$ un point.
Soient $f : X \to S$ et $g_i : Y_i \to S$, $i = 1, 2$, des familles de courbes.
Soient $c_i : X \to Y_i$ des morphismes sur $S$.
Supposons donné un isomorphisme $Y_{1, s} \cong Y_{2, s}$
des fibres, compatible avec $c_{1, s}$ et $c_{2, s}$.
Si $c_{1, s, *}\mathcal{O}_{X_s} = \mathcal{O}_{Y_{1, s}}$ et
$R^1c_{1, s, *}\mathcal{O}_{X_s} = 0$, alors il existe un
voisinage ouvert $U$ de $s$ et un isomorphisme
$Y_{1, U} \cong Y_{2, U}$ de familles de courbes sur $U$,
compatible avec l'isomorphisme donné des fibres ainsi qu'avec
$c_1$ et $c_2$.
\end{lemma}

\begin{proof}
Rappelons que $\mathcal{O}_{S, s} = \colim \mathcal{O}_S(U)$, où
la limite inductive est prise sur le système des voisinages affines $U$ de $s$.
La catégorie des espaces algébriques de présentation finie sur
l'anneau local est donc la limite inductive des catégories d'espaces algébriques
de présentation finie sur les voisinages affines de $s$.
Voir Limites d'espaces, lemme
\ref{spaces-limits-lemma-descend-finite-presentation}.
On se ramène ainsi au cas où $S$ est le spectre d'un
anneau local et où $s$ est le point fermé.

\medskip\noindent
Supposons $S = \Spec(A)$, où $A$ est un anneau local et $s$ son point fermé.
Écrivons $A = \colim A_j$, avec les $A_j$ locaux noethériens (par exemple essentiellement
de type fini sur $\mathbf{Z}$) et les homomorphismes de transition locaux.
Posons $S_j = \Spec(A_j)$, de point fermé $s_j$. On peut trouver un
$j$ et des familles de courbes $X_j \to S_j$, $Y_{j, i} \to S_j$ ;
voir le lemme \ref{lemma-curves-qs-lfp} et
Limites de champs, lemme
\ref{stacks-limits-lemma-representable-by-spaces-limit-preserving}.
Quitte à augmenter $j$, on peut trouver des morphismes
$c_{j, i} : X_j \to Y_{j, i}$ dont le changement de base à $s$ est $c_i$ ; voir
Limites d'espaces, lemme
\ref{spaces-limits-lemma-descend-finite-presentation}.
Puisque $\kappa(s) = \colim \kappa(s_j)$, on peut de même
supposer donné un isomorphisme $Y_{j, 1, s_j} \cong Y_{j, 2, s_j}$
compatible avec $c_{j, 1, s_j}$ et $c_{j, 2, s_j}$.
Enfin, les hypothèses
$c_{1, s, *}\mathcal{O}_{X_s} = \mathcal{O}_{Y_{1, s}}$ et
$R^1c_{1, s, *}\mathcal{O}_{X_s} = 0$
sont héritées par $c_{j, 1, s_j}$, car
$\{s_j \to s\}$ est un recouvrement fpqc et
$c_{1, s}$ est le changement de base de $c_{j, 1, s_j}$ par ce recouvrement
(nous omettons les détails).
Le lemme est ainsi ramené au cas traité dans le paragraphe suivant.

\medskip\noindent
Supposons que $S$ soit le spectre d'un anneau local noethérien $\Lambda$ et que
$s$ soit le point fermé. Considérons l'image schématique $Z$ de
$$
(c_1, c_2) : X \longrightarrow Y_1 \times_S Y_2
$$
L'énoncé du lemme équivaut à dire que $Z$ s'envoie
isomorphiquement sur $Y_1$ et $Y_2$ par les morphismes de projection.
Puisque la formation de l'image schématique de ce morphisme commute
au changement de base plat (Morphismes d'espaces, lemme
\ref{spaces-morphisms-lemma-flat-base-change-scheme-theoretic-image}),
nous pouvons remplacer $\Lambda$ par son complété
(Compléments d'algèbre, section \ref{more-algebra-section-permanence-completion}).

\medskip\noindent
Supposons que $S$ soit le spectre d'un anneau local noethérien complet $\Lambda$.
Observons que $X$, $Y_1$, $Y_2$ sont alors des schémas
(Compléments sur les morphismes d'espaces, lemme
\ref{spaces-more-morphisms-lemma-projective-over-complete}).
Notons $X_n$, $Y_{1, n}$, $Y_{2, n}$ les changements de base de
$X$, $Y_1$, $Y_2$ à $\Spec(\Lambda/\mathfrak m^{n + 1})$.
Rappelons que la flèche
$$
\Deformationcategory_{X_s \to Y_{2, s}} \cong
\Deformationcategory_{X_s \to Y_{1, s}} \longrightarrow
\Deformationcategory_{X_s}
$$
est une équivalence ; voir Problèmes de déformation, lemme
\ref{examples-defos-lemma-schemes-morphisms-smooth-to-base}.
Il existe donc un isomorphisme d'objets formels
$(X_n \to Y_{1, n}) \cong (X_n \to Y_{2, n})$
de $\Deformationcategory_{X_s \to Y_{1, s}}$.
Enfin, le théorème d'algébrisation de Grothendieck
(Cohomologie des schémas, lemme \ref{coherent-lemma-algebraize-morphism})
fournit un isomorphisme $Y_1 \to Y_2$ compatible avec $c_1$ et $c_2$.
\end{proof}

\begin{lemma}
\label{lemma-contract-basic}
Soit $f : X \to S$ une famille de courbes. Soit $s \in S$ un point.
Soit $h_0 : X_s \to Y_0$ un morphisme vers un schéma propre $Y_0$ sur $\kappa(s)$
tel que $h_{0, *}\mathcal{O}_{X_s} = \mathcal{O}_{Y_0}$ et
$R^1h_{0, *}\mathcal{O}_{X_s} = 0$. Alors il existe un voisinage
étale élémentaire $(U, u) \to (S, s)$, une famille de courbes $Y \to U$
et un morphisme $h : X_U \to Y$ sur $U$ dont la fibre en $u$
est isomorphe à $h_0$.
\end{lemma}

\begin{proof}
Commençons par quelques réductions ; le lecteur peut passer directement à la suite.
La question est locale sur $S$ ; on peut donc supposer $S$ affine.
Écrivons $S = \lim S_i$ comme limite projective filtrante de schémas affines $S_i$
de type fini sur $\mathbf{Z}$.
Pour un certain $i$, on peut trouver une famille de courbes $X_i \to S_i$
dont le changement de base est $X \to S$. Cela résulte du
lemme \ref{lemma-curves-qs-lfp} et de
Limites de champs, lemme
\ref{stacks-limits-lemma-representable-by-spaces-limit-preserving}.
Soit $s_i \in S_i$ l'image de $s$. Observons que
$\kappa(s) = \colim \kappa(s_i)$ et que $X_s$ est un schéma
(Espaces sur les corps, lemme
\ref{spaces-over-fields-lemma-codim-1-point-in-schematic-locus}).
Quitte à augmenter $i$, on peut supposer qu'il existe un morphisme
$h_{i, 0} : X_{i, s_i} \to Y_i$
de schémas de type fini sur $\kappa(s_i)$ dont le changement de base à
$\kappa(s)$ est $h_0$ ; voir
Limites, lemme \ref{limits-lemma-descend-finite-presentation}.
Quitte à augmenter $i$, on peut supposer $Y_i$ propre sur $\kappa(s_i)$ ; voir
Limites, lemme \ref{limits-lemma-eventually-proper}.
Soit $g_{i, 0} : Y_0 \to Y_{i, 0}$ la projection. Observons qu'il
s'agit d'un morphisme fidèlement plat, car il se déduit par changement de base de
$\Spec(\kappa(s)) \to \Spec(\kappa(s_i))$.
Par changement de base plat, on a
$$
h_{0, *}\mathcal{O}_{X_s} = g_{i, 0}^*h_{i, 0, *}\mathcal{O}_{X_{i, s_i}}
\quad\text{et}\quad
R^1h_{0, *}\mathcal{O}_{X_s} = g_{i, 0}^*Rh_{i, 0, *}\mathcal{O}_{X_{i, s_i}}
$$
voir Cohomologie des schémas, lemme
\ref{coherent-lemma-flat-base-change-cohomology}.
Par fidèle platitude, on voit que $X_i \to S_i$, $s_i \in S_i$ et
$X_{i, s_i} \to Y_i$ satisfont à toutes les hypothèses du lemme.
On est ainsi ramené au cas traité dans le paragraphe suivant.

\medskip\noindent
Supposons $S$ affine de type fini sur $\mathbf{Z}$.
Soit $\mathcal{O}_{S, s}^h$ l'hensélisé de l'anneau local
de $S$ en $s$. Observons que $\mathcal{O}_{S, s}^h$ est un anneau G d'après
Compléments d'algèbre, lemme \ref{more-algebra-lemma-henselization-G-ring} et
proposition \ref{more-algebra-proposition-ubiquity-G-ring}.
Supposons que nous puissions construire une famille de courbes
$Y' \to \Spec(\mathcal{O}_{S, s}^h)$ et un morphisme
$$
h' : X \times_S \Spec(\mathcal{O}_{S, s}^h) \longrightarrow Y'
$$
sur $\Spec(\mathcal{O}_{S, s}^h)$ dont le changement de base au point
fermé soit $h_0$. Cela suffira. En effet, utilisons d'abord l'égalité
$$
\mathcal{O}_{S, s}^h = \colim_{(U, u)} \mathcal{O}_U(U)
$$
où la limite inductive est prise sur la catégorie filtrante des
voisinages étales élémentaires (Compléments sur les morphismes, lemme
\ref{more-morphisms-lemma-describe-henselization}).
Utilisons ensuite de nouveau le fait que $Y'$ descend en une famille
$Y \to U$ pour un certain $U$ (voir les références données ci-dessus).
Puis appliquons
Limites, lemme \ref{limits-lemma-descend-finite-presentation},
pour faire descendre $h'$ en un certain $h$. On est ainsi ramené au cas
traité dans le paragraphe suivant.

\medskip\noindent
Supposons $S = \Spec(\Lambda)$, où $(\Lambda, \mathfrak m, \kappa)$
est un anneau G local noethérien hensélien et $s$ le point fermé de $S$.
Rappelons que l'application
$$
\Deformationcategory_{X_s \to Y_0} \to \Deformationcategory_{X_s}
$$
est une équivalence ; voir Problèmes de déformation, lemme
\ref{examples-defos-lemma-schemes-morphisms-smooth-to-base}.
(C'est la seule étape importante de la démonstration ; tout le reste
relève de la technique.) Notons $\Lambda^\wedge$ le complété $\mathfrak m$-adique.
Les images réciproques $X_n$ de $X$ sur $\Lambda/\mathfrak m^{n + 1}$ définissent
un objet formel $\xi$ de $\Deformationcategory_{X_s}$ sur $\Lambda^\wedge$.
L'équivalence fournit un objet formel
$\xi'$ de $\Deformationcategory_{X_s \to Y_0}$ sur $\Lambda^\wedge$.
On obtient ainsi un immense diagramme commutatif
$$
\xymatrix{
\ldots \ar[r] &
X_n \ar[r] \ar[d] &
X_{n - 1} \ar[r] \ar[d] &
\ldots \ar[r] &
X_s \ar[d] \\
\ldots \ar[r] &
Y_n \ar[r] \ar[d] &
Y_{n - 1} \ar[r] \ar[d] &
\ldots \ar[r] &
Y_0 \ar[d] \\
\ldots \ar[r] &
\Spec(\Lambda/\mathfrak m^{n + 1}) \ar[r] &
\Spec(\Lambda/\mathfrak m^n) \ar[r] &
\ldots \ar[r] &
\Spec(\kappa)
}
$$
L'objet formel $(Y_n)$ provient d'une famille de courbes
$Y' \to \Spec(\Lambda^\wedge)$ d'après
Quot, lemme \ref{quot-lemma-curves-existence}.
D'après Compléments sur les morphismes d'espaces, lemme
\ref{spaces-more-morphisms-lemma-algebraize-morphism},
on obtient un morphisme $h' : X_{\Lambda^\wedge} \to Y'$ qui induit
les morphismes donnés $X_n \to Y_n$ pour tout $n$,
et en particulier le morphisme donné $X_s \to Y_0$.

\medskip\noindent
Pour terminer, appliquons un argument standard d'algébrisation et d'approximation.
Observons d'abord que l'on peut trouver une sous-$\Lambda$-algèbre de type fini
$\Lambda \subset A \subset \Lambda^\wedge$, une famille de courbes
$Y'' \to \Spec(A)$ et un morphisme $h'' : X_A \to Y''$ sur $A$
dont le changement de base à $\Lambda^\wedge$ est $h'$.
Cela résulte de ce que $\Lambda^\wedge$ est la limite inductive filtrante
de ces anneaux $A$ et de ce que l'on peut raisonner comme précédemment en utilisant
le fait que $\Curvesstack$ est localement de présentation finie
(ce qui fournit $Y''$ sur $A$ d'après
Limites de champs, lemme
\ref{stacks-limits-lemma-representable-by-spaces-limit-preserving})
et en appliquant
Limites d'espaces, lemme \ref{spaces-limits-lemma-descend-finite-presentation},
pour faire descendre $h'$ en un certain $h''$.
On peut alors appliquer la propriété d'approximation des anneaux G
(sous la forme du théorème de Lissage des homomorphismes d'anneaux,
\ref{smoothing-theorem-approximation-property})
afin de trouver une application $A \to \Lambda$ qui induise la même application
$A \to \kappa$ que celle obtenue de $A \to \Lambda^\wedge$.
En changeant la base de $h''$ à $\Lambda$, on achève la démonstration.
\end{proof}

\begin{lemma}
\label{lemma-contract-prestable-to-stable}
Soit $f : X \to S$ une famille préstable de courbes de genre $g \geq 2$.
Il existe une factorisation $X \to Y \to S$ de $f$, où $g : Y \to S$
est une famille stable de courbes et où $c : X \to Y$ possède les propriétés
suivantes :
\begin{enumerate}
\item $\mathcal{O}_Y = c_*\mathcal{O}_X$ et $R^1c_*\mathcal{O}_X = 0$,
ces égalités restant vraies après tout changement de base par un morphisme $S' \to S$ ;
\item pour tout $s \in S$, le morphisme $c_s : X_s \to Y_s$ est la
contraction des queues et ponts rationnels étudiée dans
Courbes algébriques, section \ref{curves-section-contracting-to-stable}.
\end{enumerate}
De plus, $c : X \to Y$ est unique à isomorphisme unique près.
\end{lemma}

\begin{proof}
Soit $s \in S$. Soit $c_0 : X_s \to Y_0$ la contraction de
Courbes algébriques, section \ref{curves-section-contracting-to-stable}
(plus précisément, Courbes algébriques, lemme
\ref{curves-lemma-characterize-contraction-to-stable}).
D'après le lemme \ref{lemma-contract-basic},
il existe un voisinage étale élémentaire
$(U, u)$ et un morphisme $c : X_U \to Y$
de familles de courbes sur $U$ qui redonne
$c_0$ comme fibre en $u$.
Puisque $\omega_{Y_0}$ est ample, quitte à rétrécir $U$,
on voit que $Y \to U$ est une famille stable de genre $g$
par l'ouverture contenue dans
les lemmes \ref{lemma-stable-curves} et \ref{lemma-stable-one-piece-per-genus}.
Quitte à rétrécir $U$ une nouvelle fois, l'assertion (1) du lemme pour
$c : X_U \to Y$ résulte du lemme \ref{lemma-contract}.
De plus, la partie (2) résulte de l'unicité dans Courbes algébriques, lemme
\ref{curves-lemma-characterize-contraction-to-stable}.
On en conclut qu'un morphisme $c$ comme dans le lemme existe localement pour
la topologie étale sur $S$. Plus précisément, il existe un recouvrement étale
$\{U_i \to S\}$ et des morphismes $c_i : X_{U_i} \to Y_i$ sur $U_i$,
où $Y_i \to U_i$ est une famille stable de courbes
possédant les propriétés (1) et (2) énoncées dans le lemme.

\medskip\noindent
Pour achever la démonstration, il suffit de prouver l'unicité de $c : X \to Y$
(à isomorphisme unique près). En effet, une fois cela établi, nous
obtenons des isomorphismes
$$
\varphi_{ij} :
Y_i \times_{U_i} (U_i \times_S U_j)
\longrightarrow
Y_i \times_{U_j} (U_i \times_S U_j)
$$
satisfaisant la condition de cocycle (par unicité) sur
$U_i \times U_j \times U_k$. Puisque $\overline{\mathcal{M}_g}$
est un champ algébrique, les données de descente sont effectives,
et nous obtenons $Y \to S$. Les morphismes $c_i$ descendent en un morphisme
$c : X \to Y$ sur $S$. Enfin, les propriétés (1) et (2) pour $c$
résultent immédiatement des propriétés (1) et (2) pour $c_i$.

\medskip\noindent
Enfin, si $c_1 : X \to Y_i$, $i = 1, 2$, sont deux morphismes vers
des familles stables de courbes sur $S$ satisfaisant (1) et (2), alors
nous obtenons un morphisme $Y_1 \to Y_2$ compatible avec $c_1$ et $c_2$,
au moins localement sur $S$, grâce au lemme \ref{lemma-contract-basic-uniqueness}.
Nous omettons la vérification de l'unicité de ces morphismes
(indication : elle résulte du fait que l'image schématique
de $c_1$ est $Y_1$). Ces morphismes définis localement se recollent donc,
ce qui achève la démonstration.
\end{proof}

\begin{lemma}
\label{lemma-stabilization-morphism}
Soit $g \geq 2$. Il existe un morphisme de champs algébriques sur $\mathbf{Z}$
$$
stabilisation :
\Curvesstack^{prestable}_g
\longrightarrow
\overline{\mathcal{M}}_g
$$
qui envoie une famille préstable de courbes $X \to S$ de genre $g$
sur la famille stable $Y \to S$ qui lui est associée dans le
lemme \ref{lemma-contract-prestable-to-stable}.
\end{lemma}

\begin{proof}
Pour le voir, il suffit de vérifier que la construction du
lemme \ref{lemma-contract-prestable-to-stable}
est compatible aux changements de base (ainsi qu'aux isomorphismes, ce qui est immédiat) ;
voir l'(abus de) langage pour les champs algébriques introduit dans
Propriétés des champs, section \ref{stacks-properties-section-conventions}.
Il suffit pour cela de vérifier que les propriétés (1) et (2) du
lemme \ref{lemma-contract-prestable-to-stable} sont stables
par changement de base. C'est immédiat pour (1).
Pour (2), cela résulte soit de ce que les contractions de
Courbes algébriques, lemmes \ref{curves-lemma-contracting-rational-tails} et
\ref{curves-lemma-contracting-rational-bridges},
sont stables par extension du corps de base, soit
de ce que les conditions caractérisant les morphismes sur les fibres dans
Courbes algébriques, lemme \ref{curves-lemma-characterize-contraction-to-stable},
sont préservées par extension du corps de base.
\end{proof}



\section{Théorème de réduction stable}
\label{section-stable-reduction}

\noindent
Dans le chapitre sur la réduction semi-stable, nous avons démontré le célèbre
théorème de réduction semi-stable des courbes. Soit $K$ le corps des fractions
d'un anneau de valuation discrète $R$. Soit $C$ une courbe projective lisse
sur $K$ telle que $K = H^0(C, \mathcal{O}_C)$. Conformément à
Réduction semi-stable, définition \ref{models-definition-semistable},
on dit que $C$ a une {\it réduction semi-stable} si, au choix,
il existe une famille préstable de courbes sur $R$ de fibre générique $C$, ou
un (ce qui équivaut à tout) modèle régulier minimal de $C$ sur $R$ est préstable.
Nous montrons dans cette section que, pour les courbes de genre $g \geq 2$,
cela équivaut aussi à l'existence d'une réduction stable.

\begin{lemma}
\label{lemma-stable-reduction}
Soit $R$ un anneau de valuation discrète de corps des fractions $K$.
Soit $C$ une courbe projective lisse sur $K$ telle que
$K = H^0(C, \mathcal{O}_C)$ et de genre $g \geq 2$.
Les conditions suivantes sont équivalentes :
\begin{enumerate}
\item $C$ a une réduction semi-stable
(Réduction semi-stable, définition \ref{models-definition-semistable}) ;
\item il existe une famille stable de courbes sur $R$ de fibre générique $C$.
\end{enumerate}
\end{lemma}

\begin{proof}
Une famille stable de courbes étant aussi préstable, il est immédiat que
(2) implique (1). Réciproquement, étant donnée une famille préstable de courbes sur $R$
de fibre générique $C$, le lemme \ref{lemma-contract-prestable-to-stable}
permet de la contracter en une famille stable de courbes. Comme la fibre générique
est déjà stable, cette opération ne la modifie pas,
ce qui achève la démonstration.
\end{proof}

\noindent
Le lemme suivant affirme que la famille stable de courbes sur $R$
promise par le lemme \ref{lemma-stable-reduction}
est unique à isomorphisme unique près.

\begin{lemma}
\label{lemma-unique-stable-model}
Soit $R$ un anneau de valuation discrète de corps des fractions $K$.
Soit $C$ une courbe propre et lisse sur $K$
telle que $K = H^0(C, \mathcal{O}_C)$ et de genre $g$.
Si $X$ et $X'$ sont des modèles de $C$
(Réduction semi-stable, section \ref{models-section-models})
et si $X$ et $X'$ sont des familles stables de courbes de genre $g$ sur $R$,
alors il existe un unique isomorphisme de modèles $X \to X'$.
\end{lemma}

\begin{proof}
Soit $Y$ le modèle minimal de $C$. Rappelons que $Y$ existe, est
unique et est au plus nodal de dimension relative $1$ sur $R$ ; voir
Réduction semi-stable,
proposition \ref{models-proposition-exists-minimal-model} et
lemmes \ref{models-lemma-minimal-model-unique} et
\ref{models-lemma-semistable} (ce dernier s'applique puisque l'on dispose de $X$).
Il existe un morphisme de contraction
$$
Y \longrightarrow Z
$$
tel que $Z$ soit une famille stable de courbes de genre $g$ sur $R$
(lemme \ref{lemma-contract-prestable-to-stable}). Nous affirmons
qu'il existe un unique isomorphisme de modèles $X \to Z$.
Par symétrie, il en va de même pour $X'$, ce qui achèvera la démonstration.

\medskip\noindent
D'après Réduction semi-stable, lemme \ref{models-lemma-blowup-at-worst-nodal},
il existe une suite
$$
X_m \to \ldots \to X_1 \to X_0 = X
$$
telle que $X_{i + 1} \to X_i$ soit l'éclatement d'un point fermé
$x_i$ où $X_i$ est singulier, que $X_i \to \Spec(R)$ soit au plus nodal
de dimension relative $1$, et que $X_m$ soit régulier.
D'après Réduction semi-stable, lemme \ref{models-lemma-pre-exists-minimal-model},
il existe une suite
$$
X_m = Y_n \to Y_{n - 1} \to \ldots \to Y_1 \to Y_0 = Y
$$
de modèles propres réguliers de $C$, dont chaque morphisme est une
contraction d'une courbe exceptionnelle de première espèce\footnote{En fait,
on a $X_m = Y$, c'est-à-dire que $X_m$ ne contient aucune courbe exceptionnelle
de première espèce. Nous invitons le lecteur à le vérifier,
car cela simplifie quelque peu la démonstration.}.
D'après Réduction semi-stable, lemme \ref{models-lemma-blowdown-at-worst-nodal},
chaque $Y_i$ est au plus nodal de dimension relative $1$ sur $R$.
Pour établir notre affirmation, il suffit de montrer qu'il existe un isomorphisme
$X \to Z$ compatible avec les morphismes $X_m \to X$
et $X_m = Y_n \to Y \to Z$. Soit $s \in \Spec(R)$ le point fermé.
En appliquant soit le
lemme \ref{lemma-contract-basic-uniqueness}, soit le
lemme \ref{lemma-contract-prestable-to-stable},
on se ramène à démontrer que les morphismes
$X_{m, s} \to X_s$ et
$X_{m, s} \to Z_s$
sont tous deux égaux au morphisme canonique de
Courbes algébriques, lemme \ref{curves-lemma-characterize-contraction-to-stable}.

\medskip\noindent
Pour un morphisme $c : U \to V$ de schémas sur $\kappa(s)$,
on dit que $c$ possède la propriété (*) si $\dim(U_v) \leq 1$ pour $v \in V$,
$\mathcal{O}_V = c_*\mathcal{O}_U$ et $R^1c_*\mathcal{O}_U = 0$.
Cette propriété est stable par composition.
Puisque $X_s$ et $Z_s$ sont tous deux des courbes stables de genre $g$ sur $\kappa(s)$,
il suffit de montrer que chacun des morphismes $Y_s \to Z_s$,
$X_{i + 1, s} \to X_{i, s}$ et $Y_{i + 1, s} \to Y_{i, s}$
possède la propriété (*) ; voir
Courbes algébriques, lemme \ref{curves-lemma-characterize-contraction-to-stable}.

\medskip\noindent
La propriété (*) est satisfaite par $Y_s \to Z_s$ par construction.

\medskip\noindent
Les morphismes $c : X_{i + 1, s} \to X_{i, s}$ sont construits et étudiés
dans la démonstration de
Réduction semi-stable, lemme \ref{models-lemma-blowup-at-worst-nodal}.
Il suffit de vérifier (*) localement pour la topologie étale sur $X_{i, s}$.
Il suffit donc de vérifier (*) pour le changement de base du morphisme
« $X_1 \to X_0$ » de Réduction semi-stable, exemple \ref{models-example-blowup},
à $R/\pi R$.
Nous laissons le calcul explicite au lecteur.

\medskip\noindent
Le morphisme $c : Y_{i + 1, s} \to Y_{i, s}$ est la restriction
de la contraction d'une courbe exceptionnelle $E \subset Y_{i + 1}$
de première espèce ; autrement dit, $b : Y_{i + 1} \to Y_i$ est une contraction de $E$,
c'est-à-dire que $b$ est l'éclatement d'un point régulier de la surface $Y_i$
(Résolution des surfaces, section \ref{resolve-section-minus-one}).
On a alors $\mathcal{O}_{Y_i} = b_*\mathcal{O}_{Y_{i + 1}}$ et
$R^1b_*\mathcal{O}_{Y_{i + 1}} = 0$ ; voir par exemple
Résolution des surfaces, lemme
\ref{resolve-lemma-cohomology-of-blowup}.
Nous en concluons que $\mathcal{O}_{Y_{i, s}} = c_*\mathcal{O}_{Y_{i + 1, s}}$
et $R^1c_*\mathcal{O}_{Y_{i + 1, s}} = 0$ d'après
Compléments sur les morphismes, lemmes \ref{more-morphisms-lemma-check-h1-fibre-zero},
\ref{more-morphisms-lemma-h1-fibre-zero} et
\ref{more-morphisms-lemma-h1-fibre-zero-check-h0-kappa}
(ce dernier ne donne que la surjectivité de
$\mathcal{O}_{Y_{i, s}} \to c_*\mathcal{O}_{Y_{i + 1, s}}$, mais
l'injectivité résulte aisément de ce que $Y_{i, s}$ est réduit
et que $c$ ne modifie la situation qu'au-dessus d'un point fermé).
Cela achève la démonstration.
\end{proof}

\noindent
Du lemme \ref{lemma-stable-reduction} et du
théorème de Réduction semi-stable \ref{models-theorem-semistable-reduction},
on déduit immédiatement le théorème de réduction stable.

\begin{theorem}
\label{theorem-stable-reduction}
\begin{reference}
\cite[corollaire 2.7]{DM}
\end{reference}
Soit $R$ un anneau de valuation discrète de corps des fractions $K$. Soit $C$ une
courbe projective lisse sur $K$ telle que $H^0(C, \mathcal{O}_C) = K$
et de genre $g \geq 2$. Alors :
\begin{enumerate}
\item il existe une extension d'anneaux de valuation discrète $R \subset R'$
induisant une extension finie séparable des corps des fractions $K'/K$, et
une famille stable de courbes $Y \to \Spec(R')$ de genre $g$ telle que
$Y_{K'} \cong C_{K'}$ sur $K'$ ;
\item il existe une extension finie séparable $L/K$ et une famille stable
de courbes $Y \to \Spec(A)$ de genre $g$, où $A \subset L$
est la clôture intégrale de $R$ dans $L$, telles que
$Y_L \cong C_L$ sur $L$.
\end{enumerate}
\end{theorem}

\begin{proof}
La partie (1) résulte immédiatement du lemme \ref{lemma-stable-reduction} et du
théorème de Réduction semi-stable \ref{models-theorem-semistable-reduction}.

\medskip\noindent
Démontrons (2). Soit $L/K$ l'extension finie séparable obtenue dans la partie (3) du
théorème de Réduction semi-stable \ref{models-theorem-semistable-reduction}.
Soit $A \subset L$ la clôture intégrale de $R$.
Rappelons que $A$ est un anneau de Dedekind fini sur $R$, possédant
un nombre fini d'idéaux maximaux $\mathfrak m_1, \ldots, \mathfrak m_n$ ; voir
Compléments d'algèbre, remarque \ref{more-algebra-remark-finite-separable-extension}.
Posons $S = \Spec(A)$, $S_i = \Spec(A_{\mathfrak m_i})$,
$U = \Spec(L)$ et $U_i = S_i \setminus \{\mathfrak m_i\}$.
Observons que $U \cong U_i$ pour $i = 1, \ldots, n$.
Posons $X = C_L$, considéré comme schéma sur le sous-schéma ouvert $U$ de $S$.
Par le choix de $L$ et de $A$ et par le lemme \ref{lemma-stable-reduction},
nous disposons de familles stables de courbes $X_i \to S_i$ et d'isomorphismes
$X \times_U U_i \cong X_i \times_{S_i} U_i$.
D'après Limites d'espaces, lemme
\ref{spaces-limits-lemma-glueing-near-multiple-closed-points},
on peut trouver un morphisme de présentation finie $Y \to S$
dont le changement de base à $S_i$ est isomorphe à $X_i$ pour $i = 1, \ldots, n$.
On peut aussi utiliser le fait que $S = \bigcup_{i = 1, \ldots, n} S_i$
est un recouvrement ouvert de $S$ et que $S_i \cap S_j = U$ pour $i \not = j$,
puis appliquer $n - 1$ fois
Limites d'espaces, lemme \ref{spaces-limits-lemma-relative-glueing},
afin d'obtenir $Y \to S$ dont le
changement de base à $S_i$ est isomorphe à $X_i$ pour $i = 1, \ldots, n$.
Il est clair que $Y \to S$ est la famille stable de courbes recherchée.
\end{proof}






\section{Propriétés du champ des courbes stables}
\label{section-properties-stable}

\noindent
Dans cette section, nous démontrons le résultat structurel fondamental pour
$\overline{\mathcal{M}}_g$ lorsque $g \geq 2$.

\begin{lemma}
\label{lemma-stable-separated}
Soit $g \geq 2$. Le champ $\overline{\mathcal{M}}_g$ est séparé.
\end{lemma}

\begin{proof}
Cet énoncé signifie que le morphisme
$\overline{\mathcal{M}}_g \to \Spec(\mathbf{Z})$ est séparé.
Nous allons le démontrer au moyen du critère valuatif noethérien raffiné
énoncé dans
Compléments sur les morphismes de champs, lemme
\ref{stacks-more-morphisms-lemma-refined-valuative-criterion-separated}.

\medskip\noindent
Puisque $\overline{\mathcal{M}}_g$ est un sous-champ ouvert de
$\Curvesstack$, le morphisme $\overline{\mathcal{M}}_g \to \Spec(\mathbf{Z})$
est quasi-séparé et
localement de présentation finie d'après le lemme \ref{lemma-curves-qs-lfp}.
En particulier, le champ $\overline{\mathcal{M}}_g$ est localement
noethérien (Morphismes de champs, lemme
\ref{stacks-morphisms-lemma-locally-finite-type-locally-noetherian}).
D'après le lemme \ref{lemma-smooth-dense-in-stable}, l'immersion ouverte
$\mathcal{M}_g \to \overline{\mathcal{M}}_g$
est d'image dense. De plus, $\mathcal{M}_g \to \overline{\mathcal{M}}_g$
est quasi-compact (Morphismes de champs, lemme
\ref{stacks-morphisms-lemma-locally-closed-in-noetherian}),
donc de type fini. Toutes les hypothèses préliminaires de
Compléments sur les morphismes de champs, lemme
\ref{stacks-more-morphisms-lemma-refined-valuative-criterion-separated},
sont ainsi satisfaites pour les morphismes
$$
\mathcal{M}_g \to \overline{\mathcal{M}}_g
\quad\text{et}\quad
\overline{\mathcal{M}}_g \to \Spec(\mathbf{Z})
$$
et il suffit de vérifier ce qui suit : pour tout diagramme $2$-commutatif
$$
\xymatrix{
\Spec(K) \ar[r] \ar[d] &
\mathcal{M}_g \ar[r] &
\overline{\mathcal{M}}_g \ar[d] \\
\Spec(R) \ar[rr] \ar@{..>}[rru] & & \Spec(\mathbf{Z})
}
$$
où $R$ est un anneau de valuation discrète de corps des fractions $K$,
la catégorie des flèches pointillées est soit vide, soit un sétoïde ayant exactement
une classe d'isomorphisme. (Observons qu'il n'est pas nécessaire de trop se soucier
des $2$-flèches ; voir Morphismes de champs, lemme
\ref{stacks-morphisms-lemma-cat-dotted-arrows-independent}).
En explicitant cet énoncé et en utilisant le fait que
$\mathcal{M}_g$, resp.\ $\overline{\mathcal{M}}_g$, sont
les champs algébriques qui paramètrent les familles lisses, resp.\ stables,
de courbes de genre $g$, on constate que l'assertion à démontrer
est exactement le résultat d'unicité énoncé et démontré dans le
lemme \ref{lemma-unique-stable-model}.
\end{proof}

\begin{lemma}
\label{lemma-stable-quasi-compact}
Soit $g \geq 2$. Le champ $\overline{\mathcal{M}}_g$ est quasi-compact.
\end{lemma}

\begin{proof}
Nous utiliserons les notations de la section \ref{section-polarized-curves}.
Considérons le sous-ensemble
$$
T \subset |\textit{PolarizedCurves}|
$$
des points $\xi$ pour lesquels il existe un corps $k$ et une paire
$(X, \mathcal{L})$ sur $k$ représentant $\xi$
et possédant les deux propriétés suivantes :
\begin{enumerate}
\item $X$ est une courbe stable de genre $g$ ;
\item $\mathcal{L} = \omega_X^{\otimes 3}$.
\end{enumerate}
Manifestement, par l'application continue
$$
|\textit{PolarizedCurves}|
\longrightarrow
|\Curvesstack|
$$
l'image de l'ensemble $T$ est exactement le sous-ensemble ouvert
$$
|\overline{\mathcal{M}}_g| \subset |\Curvesstack|
$$
Il suffit donc de montrer que $T$ est quasi-compact.
D'après le lemme \ref{lemma-polarized-curves-in-polarized}, on a une
immersion ouverte et fermée
$$
|\textit{PolarizedCurves}| \subset |\Polarizedstack|
$$
Il suffit donc de démontrer la quasi-compacité de $T$ comme sous-ensemble de
$|\Polarizedstack|$. Nous appliquons pour cela le critère de
Champs de modules, lemme \ref{moduli-lemma-bounded-polarized}.
Observons d'abord que, pour $(X, \mathcal{L})$
comme ci-dessus, le polynôme de Hilbert $P$ est la fonction
$P(t) = (6g - 6)t + (1 - g)$ d'après Riemann--Roch ; voir
Courbes algébriques, lemme \ref{curves-lemma-rr}.
Observons ensuite que $H^1(X, \mathcal{L}) = 0$
et que $\mathcal{L}$ est très ample d'après
Courbes algébriques, lemme \ref{curves-lemma-tricanonical}.
Cela signifie exactement que, pour $n = P(3) - 1$,
il existe une immersion fermée
$$
i : X \longrightarrow \mathbf{P}^n_k
$$
telle que $\mathcal{L} = i^*\mathcal{O}_{\mathbf{P}^1_k}(1)$,
comme souhaité.
\end{proof}

\noindent
Voici le théorème principal de cette section.

\begin{theorem}
\label{theorem-stable-smooth-proper}
Soit $g \geq 2$. Le champ algébrique $\overline{\mathcal{M}}_g$ est un
champ de Deligne-Mumford, propre et lisse sur $\Spec(\mathbf{Z})$.
De plus, le lieu $\mathcal{M}_g$ qui paramètre les courbes lisses
est un sous-champ ouvert dense.
\end{theorem}

\begin{proof}
La plupart des propriétés mentionnées dans l'énoncé ont déjà été démontrées.
La lissité résulte du lemme \ref{lemma-stable-curves-smooth}.
La propriété de Deligne-Mumford résulte du lemme \ref{lemma-stable-curves-deligne-mumford}.
L'ouverture de $\mathcal{M}_g$ résulte du lemme \ref{lemma-smooth-dense-in-stable}.
Nous savons que $\overline{\mathcal{M}}_g \to \Spec(\mathbf{Z})$
est séparé d'après le lemme \ref{lemma-stable-separated} et que
$\overline{\mathcal{M}}_g$ est quasi-compact d'après le
lemme \ref{lemma-stable-quasi-compact}.
Ainsi, pour montrer que $\overline{\mathcal{M}}_g \to \Spec(\mathbf{Z})$
est propre et achever la démonstration, on peut appliquer
Compléments sur les morphismes de champs, lemme
\ref{stacks-more-morphisms-lemma-refined-valuative-criterion-proper},
aux morphismes $\mathcal{M}_g \to \overline{\mathcal{M}}_g$ et
$\overline{\mathcal{M}}_g \to \Spec(\mathbf{Z})$.
Il suffit donc de vérifier ce qui suit : pour tout diagramme $2$-commutatif
$$
\xymatrix{
\Spec(K) \ar[r] \ar[d]_j &
\mathcal{M}_g \ar[r] &
\overline{\mathcal{M}}_g \ar[d] \\
\Spec(A) \ar[rr] & & \Spec(\mathbf{Z})
}
$$
où $A$ est un anneau de valuation discrète de corps des fractions $K$, il
existe une extension de corps $K'/K$ et un anneau de valuation $A' \subset K'$
dominant $A$, tels que la catégorie des flèches pointillées du
diagramme induit
$$
\xymatrix{
\Spec(K') \ar[r] \ar[d]_{j'} & \overline{\mathcal{M}}_g \ar[d] \\
\Spec(A') \ar[r] \ar@{..>}[ru] & \Spec(\mathbf{Z})
}
$$
soit non vide (Morphismes de champs, définition
\ref{stacks-morphisms-definition-fill-in-diagram}).
(Observons qu'il n'est pas nécessaire de trop se soucier des
$2$-flèches ; voir Morphismes de champs, lemme
\ref{stacks-morphisms-lemma-cat-dotted-arrows-independent}).
En explicitant cet énoncé et en utilisant le fait que
$\mathcal{M}_g$, resp.\ $\overline{\mathcal{M}}_g$, sont les champs algébriques
qui paramètrent les familles lisses, resp.\ stables, de courbes de genre $g$,
on constate que l'assertion à démontrer est exactement le résultat contenu
dans le théorème de réduction stable, c'est-à-dire le théorème
\ref{theorem-stable-reduction}.
\end{proof}










\begin{multicols}{2}[\section{Autres chapitres}]
\noindent
Préliminaires
\begin{enumerate}
\item \hyperref[introduction-section-phantom]{Introduction}
\item \hyperref[conventions-section-phantom]{Conventions}
\item \hyperref[sets-section-phantom]{Théorie des ensembles}
\item \hyperref[categories-section-phantom]{Catégories}
\item \hyperref[topology-section-phantom]{Topologie}
\item \hyperref[sheaves-section-phantom]{Faisceaux sur les espaces}
\item \hyperref[sites-section-phantom]{Sites et faisceaux}
\item \hyperref[stacks-section-phantom]{Champs}
\item \hyperref[fields-section-phantom]{Corps}
\item \hyperref[algebra-section-phantom]{Algèbre commutative}
\item \hyperref[brauer-section-phantom]{Groupes de Brauer}
\item \hyperref[homology-section-phantom]{Algèbre homologique}
\item \hyperref[derived-section-phantom]{Catégories dérivées}
\item \hyperref[simplicial-section-phantom]{Méthodes simpliciales}
\item \hyperref[more-algebra-section-phantom]{Compléments d'algèbre}
\item \hyperref[smoothing-section-phantom]{Lissage des morphismes d'anneaux}
\item \hyperref[modules-section-phantom]{Faisceaux de modules}
\item \hyperref[sites-modules-section-phantom]{Modules sur les sites}
\item \hyperref[injectives-section-phantom]{Objets injectifs}
\item \hyperref[cohomology-section-phantom]{Cohomologie des faisceaux}
\item \hyperref[sites-cohomology-section-phantom]{Cohomologie sur les sites}
\item \hyperref[dga-section-phantom]{Algèbre différentielle graduée}
\item \hyperref[dpa-section-phantom]{Algèbre à puissances divisées}
\item \hyperref[sdga-section-phantom]{Faisceaux différentiels gradués}
\item \hyperref[hypercovering-section-phantom]{Hyperrecouvrements}
\end{enumerate}
Schémas
\begin{enumerate}
\setcounter{enumi}{25}
\item \hyperref[schemes-section-phantom]{Schémas}
\item \hyperref[constructions-section-phantom]{Constructions de schémas}
\item \hyperref[properties-section-phantom]{Propriétés des schémas}
\item \hyperref[morphisms-section-phantom]{Morphismes de schémas}
\item \hyperref[coherent-section-phantom]{Cohomologie des schémas}
\item \hyperref[divisors-section-phantom]{Diviseurs}
\item \hyperref[limits-section-phantom]{Limites de schémas}
\item \hyperref[varieties-section-phantom]{Variétés}
\item \hyperref[topologies-section-phantom]{Topologies sur les schémas}
\item \hyperref[descent-section-phantom]{Descente}
\item \hyperref[perfect-section-phantom]{Catégories dérivées de schémas}
\item \hyperref[more-morphisms-section-phantom]{Compléments sur les morphismes}
\item \hyperref[flat-section-phantom]{Compléments sur la platitude}
\item \hyperref[groupoids-section-phantom]{Schémas en groupoïdes}
\item \hyperref[more-groupoids-section-phantom]{Compléments sur les schémas en groupoïdes}
\item \hyperref[etale-section-phantom]{Morphismes étales de schémas}
\end{enumerate}
Thèmes de théorie des schémas
\begin{enumerate}
\setcounter{enumi}{41}
\item \hyperref[chow-section-phantom]{Homologie de Chow}
\item \hyperref[intersection-section-phantom]{Théorie de l'intersection}
\item \hyperref[pic-section-phantom]{Schémas de Picard des courbes}
\item \hyperref[weil-section-phantom]{Théories cohomologiques de Weil}
\item \hyperref[adequate-section-phantom]{Modules adéquats}
\item \hyperref[dualizing-section-phantom]{Complexes dualisants}
\item \hyperref[duality-section-phantom]{Dualité pour les schémas}
\item \hyperref[discriminant-section-phantom]{Discriminants et différentes}
\item \hyperref[derham-section-phantom]{Cohomologie de de Rham}
\item \hyperref[local-cohomology-section-phantom]{Cohomologie locale}
\item \hyperref[algebraization-section-phantom]{Géométrie algébrique et formelle}
\item \hyperref[curves-section-phantom]{Courbes algébriques}
\item \hyperref[resolve-section-phantom]{Résolution des surfaces}
\item \hyperref[models-section-phantom]{Réduction semi-stable}
\item \hyperref[functors-section-phantom]{Foncteurs et morphismes}
\item \hyperref[equiv-section-phantom]{Catégories dérivées de variétés}
\item \hyperref[pione-section-phantom]{Groupes fondamentaux de schémas}
\item \hyperref[etale-cohomology-section-phantom]{Cohomologie étale}
\item \hyperref[crystalline-section-phantom]{Cohomologie cristalline}
\item \hyperref[proetale-section-phantom]{Cohomologie pro-étale}
\item \hyperref[relative-cycles-section-phantom]{Cycles relatifs}
\item \hyperref[more-etale-section-phantom]{Compléments de cohomologie étale}
\item \hyperref[trace-section-phantom]{La formule des traces}
\end{enumerate}
Espaces algébriques
\begin{enumerate}
\setcounter{enumi}{64}
\item \hyperref[spaces-section-phantom]{Espaces algébriques}
\item \hyperref[spaces-properties-section-phantom]{Propriétés des espaces algébriques}
\item \hyperref[spaces-morphisms-section-phantom]{Morphismes d'espaces algébriques}
\item \hyperref[decent-spaces-section-phantom]{Espaces algébriques décents}
\item \hyperref[spaces-cohomology-section-phantom]{Cohomologie des espaces algébriques}
\item \hyperref[spaces-limits-section-phantom]{Limites d'espaces algébriques}
\item \hyperref[spaces-divisors-section-phantom]{Diviseurs sur les espaces algébriques}
\item \hyperref[spaces-over-fields-section-phantom]{Espaces algébriques sur des corps}
\item \hyperref[spaces-topologies-section-phantom]{Topologies sur les espaces algébriques}
\item \hyperref[spaces-descent-section-phantom]{Descente et espaces algébriques}
\item \hyperref[spaces-perfect-section-phantom]{Catégories dérivées d'espaces}
\item \hyperref[spaces-more-morphisms-section-phantom]{Compléments sur les morphismes d'espaces}
\item \hyperref[spaces-flat-section-phantom]{Platitude sur les espaces algébriques}
\item \hyperref[spaces-groupoids-section-phantom]{Groupoïdes dans les espaces algébriques}
\item \hyperref[spaces-more-groupoids-section-phantom]{Compléments sur les groupoïdes dans les espaces}
\item \hyperref[bootstrap-section-phantom]{Amorçage}
\item \hyperref[spaces-pushouts-section-phantom]{Sommes amalgamées d'espaces algébriques}
\end{enumerate}
Thèmes de géométrie
\begin{enumerate}
\setcounter{enumi}{81}
\item \hyperref[spaces-chow-section-phantom]{Groupes de Chow des espaces}
\item \hyperref[groupoids-quotients-section-phantom]{Quotients de groupoïdes}
\item \hyperref[spaces-more-cohomology-section-phantom]{Compléments sur la cohomologie des espaces}
\item \hyperref[spaces-simplicial-section-phantom]{Espaces simpliciaux}
\item \hyperref[spaces-duality-section-phantom]{Dualité pour les espaces}
\item \hyperref[formal-spaces-section-phantom]{Espaces algébriques formels}
\item \hyperref[restricted-section-phantom]{Algébrisation des espaces formels}
\item \hyperref[spaces-resolve-section-phantom]{Résolution des surfaces revisitée}
\end{enumerate}
Théorie des déformations
\begin{enumerate}
\setcounter{enumi}{89}
\item \hyperref[formal-defos-section-phantom]{Théorie formelle des déformations}
\item \hyperref[defos-section-phantom]{Théorie des déformations}
\item \hyperref[cotangent-section-phantom]{Le complexe cotangent}
\item \hyperref[examples-defos-section-phantom]{Problèmes de déformation}
\end{enumerate}
Champs algébriques
\begin{enumerate}
\setcounter{enumi}{93}
\item \hyperref[algebraic-section-phantom]{Champs algébriques}
\item \hyperref[examples-stacks-section-phantom]{Exemples de champs}
\item \hyperref[stacks-sheaves-section-phantom]{Faisceaux sur les champs algébriques}
\item \hyperref[criteria-section-phantom]{Critères de représentabilité}
\item \hyperref[artin-section-phantom]{Axiomes d'Artin}
\item \hyperref[quot-section-phantom]{Espaces de Quot et de Hilbert}
\item \hyperref[stacks-properties-section-phantom]{Propriétés des champs algébriques}
\item \hyperref[stacks-morphisms-section-phantom]{Morphismes de champs algébriques}
\item \hyperref[stacks-limits-section-phantom]{Limites de champs algébriques}
\item \hyperref[stacks-cohomology-section-phantom]{Cohomologie des champs algébriques}
\item \hyperref[stacks-perfect-section-phantom]{Catégories dérivées de champs}
\item \hyperref[stacks-introduction-section-phantom]{Introduction aux champs algébriques}
\item \hyperref[stacks-more-morphisms-section-phantom]{Compléments sur les morphismes de champs}
\item \hyperref[stacks-geometry-section-phantom]{La géométrie des champs}
\end{enumerate}
Thèmes de théorie des espaces de modules
\begin{enumerate}
\setcounter{enumi}{107}
\item \hyperref[moduli-section-phantom]{Champs de modules}
\item \hyperref[moduli-curves-section-phantom]{Espaces de modules de courbes}
\end{enumerate}
Divers
\begin{enumerate}
\setcounter{enumi}{109}
\item \hyperref[examples-section-phantom]{Exemples}
\item \hyperref[exercises-section-phantom]{Exercices}
\item \hyperref[guide-section-phantom]{Guide bibliographique}
\item \hyperref[desirables-section-phantom]{Desiderata}
\item \hyperref[coding-section-phantom]{Style de programmation}
\item \hyperref[obsolete-section-phantom]{Obsolète}
\item \hyperref[fdl-section-phantom]{Licence de documentation libre GNU}
\item \hyperref[index-section-phantom]{Index généré automatiquement}
\end{enumerate}
\end{multicols}


\bibliography{my}
\bibliographystyle{amsalpha}

\end{document}
